% Nichidai_MasterExam
\documentclass[dvipdfmx]{jsreport}
\usepackage{soupreamble}


\title{日本大学理工学部理工学研究科数学専攻 過去問の問題と解答(非公式)}
\date{} %日付 とりあえず空白で
\begin{document}
\maketitle %タイトル出力

\begin{abstract} %注意書き的なアブスト
	本PDFは日本大学理工学部理工学研究科数学専攻の博士前期課程(修士)の一般試験の問題・解答である。\\
	解答作成は教員らではなく，有志で行なったため必ずしも正解とは限らない。また\TeX 化に伴い一部問題文の意図が変わらない程度に文字・記号の変更を行なった部分がある。
\end{abstract}

\tableofcontents %目次出力
%年度<chap>-> 一般第○期(1or 2or 3)<sec>-> 第△問<subsec>
%正しい問題文ラベルは%%Ra_b(st/nd/rd)_c


\chapter{令和4年度}
\section{一般第3期}
\begin{tcolorbox} %%R4_3rd_1
	$\fbox{1}$ $V, W$を実数上の線形空間とする。
\begin{enumerate}
	\item 写像$T\colon V\to W$が線形写像であることの定義を書け。
	\item $T\colon V\to W$を線形写像とし、$v_1, \ldots, v_n\in V$とする。このとき$T(v_1), \ldots, T(v_n)$が線形独立ならば、$v_1, \ldots, v_n$が線形独立であることを示せ。
	\item 次で定義される$S, T$が線形写像か否か理由をつけて判定せよ。
\begin{enumerate_alpha}
	\item $S\colon \mathcal{C}\to \mathcal{C},\ \ (S(f))(x)= f(\sqrt{3})+ f(x^2)$
	\item $T\colon \mathcal{C}\to \mathcal{C},\ \ (T(f))(x)= f(x)^2$
\end{enumerate_alpha}
\end{enumerate}
\end{tcolorbox}


\begin{souanswer} %%R4_3rd_1_ans
	$\fbox{1}$
\begin{enumerate}
	\item 任意のベクトル$\bm{u}, \bm{v}\in V$とスカラー$c\in \R$に対して
\begin{enumerate_roman}
	\item\label{law:加法性} $T(\bm{u}+ \bm{v})= T(\bm{u})+ T(\bm{v})$
	\item\label{law:斉次性} $T(c\bm{v})= cT(\bm{v})$
\end{enumerate_roman}
	を満たす。
	\item 線型独立の定義より、スカラー$c_1, c_2, \ldots, c_n\in \R$に対して、以下の線型関係を仮定する。
\begin{equation*}
	a_1v_1+ a_2v_2+ \cdots+ a_nv_n= \bm{0}_V
\end{equation*}
	この両辺に$T$を作用させると、
\begin{equation*}
	T(a_1v_1+ a_2v_2+ \cdots+ a_nv_n)= T(\bm{0}_V)
\end{equation*}
	ここで$T$は線型写像であったから性質\ref{law:加法性}, \ref{law:斉次性}より
\begin{equation*}
	T(a_1v_1)+ T(a_2v_2)+ \cdots+ T(a_nv_n)= T(\bm{0}_V)
\end{equation*}
\begin{equation*}
	a_1T(v_1)+ a_2T(v_2)+ \cdots+ a_nT(v_n)= \bm{0}_W
\end{equation*}
	また仮定より$T(v_1), \ldots, T(v_n)$は線型独立であったからこの式が成り立つのは
\begin{equation*}
	a_1= a_2= \cdots= a_n= 0
\end{equation*}
	のときのみ。
	\item いずれも性質\ref{law:加法性}, \ref{law:斉次性}を満たすか確認すればよい。
\begin{enumerate_alpha}
	\item $S$は線型写像である。任意の関数$f, g\in \mathcal{C}$とスカラー$c\in \R$を用いて示す。
\begin{enumerate_roman}
	\item 
\begin{align*}
	(S(f+ g))(x)
	&= (f+ g)(\sqrt{3})+ (f+ g)(x^2)\\
	&= f(\sqrt{3})+ g(\sqrt{3})+ f(x^2)+ g(x^2)\\
	&= f(\sqrt{3})+ f(x^2)+ g(\sqrt{3})+ g(x^2)\\
	&= (S(f))(x)+ (S(g))(x)
\end{align*}
	\item 
\begin{align*}
	(S(cf))(x)
	&= (cf)(\sqrt{3})+ (cf)(x^2)\\
	&= c(f(\sqrt{3})+ f(x^2))\\
	&= c(S(f))(x)
\end{align*}
\end{enumerate_roman}
	\item $T$は線型写像ではない。\ref{law:斉次性}についての反例を挙げる。
	$f(x)= x, c= 2$とする。
\begin{align*}
	(T(cf))(x)&= (2f(x))^2= (2x)^2= 4x^2\\
	f(cx)^2&= 2f(x)^2= 2x^2
\end{align*}
	\ref{law:斉次性}を満たさないため、$T$は線型写像ではない。
\end{enumerate_alpha}
\end{enumerate}
\end{souanswer}

\begin{tcolorbox} %%R4_3rd_2
	$\fbox{2}$ $V$を複素数上の線形空間とし、$\innerproduct{\bullet}{\bullet}$で定められる内積を持つとする。複素数$z$の複素共役を$\bar{z}$、虚数単位を$i$とする。

\begin{enumerate}
	\item $W$を$V$の部分空間としたとき、$W$の直交補空間の定義を書け。
	\item $v_1, \cdots, v_n\in V$が正規直交系であるとする。また$v\in V$が複素数$c_1, \ldots, c_n$を用いて$v= \sum_{j= 1}^n c_j v_j$とかけるとする。このとき$c_j= \innerproduct{v}{v_j}$であることを示せ。
	\item $V= \C^3$とし、$\bm{a}= \begin{pmatrix}a_1\\ a_2\\ a_3\end{pmatrix},\ \bm{b}= \begin{pmatrix}b_1\\ b_2\\ b_3\end{pmatrix}$に対して、$\innerproduct{\bm{a}}{\bm{b}}= \sum_{j= 1}^3 a_j\bar{b_j}$で内積を定義する。このとき$\begin{pmatrix}1\\ 2i\\ 3\end{pmatrix},\ \begin{pmatrix}i\\ -1\\ 0\end{pmatrix}$が生成する部分空間の直交基底を1つ求めなさい。(正規直交基底である必要はない。)
\end{enumerate}
\end{tcolorbox}

\begin{souanswer} %%R4_3rd_2_ans
	$\fbox{2}$
\begin{enumerate}
	\item $V$を内積空間とし、$W\subset V$をその部分空間とする。このとき$W$のすべてのベクトルと直交するようなベクトル全体の集合
\begin{equation*}
	V^\perp= \{x\in V; \innerproduct{x}{a}= 0, \forall a\in V\}
\end{equation*}

	\item 定義式より
\begin{equation*}
	c_j= \innerproduct{\sum_{j= 1}^n c_iv_i}{v_j}
\end{equation*}
	ここで内積の性質より和の$\sum$と係数を外に出して、
\begin{equation*}
	c_j= \sum_{j= 1}^n c_i\innerproduct{v_i}{v_j}
\end{equation*}
	また$\{v_1, \ldots, v_j\}$は正規直交基底であるからKroneckerのdeltaを用いて
\begin{equation*}
	\innerproduct{v_i}{v_j}= \delta_{ij}= 
	\begin{cases}
		0\ \ (i\ne j)\\
		1\ \ (i= j)
	\end{cases}
\end{equation*}
	とかける。よって
\begin{align*}
	c_j
	&= \innerproduct{v}{v_j}\\
	&= c_1\cdot 0+ \cdots+ c_j\cdot 1+ \cdots+ c_n\cdot 0\\
	&= c_j	
\end{align*}

	\item $u_1= \begin{pmatrix}1\\ 2i\\ 3\end{pmatrix}, u_2= \begin{pmatrix}i\\ -1\\ 0\end{pmatrix}$とする。\\
	$u_1= v_1= \begin{pmatrix}1\\ 2i\\ 3\end{pmatrix}$。$u_2$から$v_1$方向の射影を取り除く。
\begin{equation*}
	v_2= u_2- \frac{\innerproduct{u_2}{v_1}}{\norm{v_1}^2}v_1
\end{equation*}
\begin{itemize}
	\item 内積$\innerproduct{u_2}{v_1}$を計算する。
\begin{equation*}
	\begin{pmatrix}1\\ 2i\\ 3\end{pmatrix}\begin{pmatrix}i\\ -1\\ 0\end{pmatrix}= i+ 2i- 0= 3i
\end{equation*}

	\item ノルム$\norm{v_1}^2$を計算する。
\begin{equation*}
	\norm{v_1}^2= \sqrt{1^2+ (2i)^2+ 3^2}^2= \sqrt{1+ 4+ 9}^2= 14
\end{equation*}
\end{itemize}
	よって
\begin{align*}
	v_2
	&= u_2- \frac{\innerproduct{u_2}{v_1}}{\norm{v_1}^2}v_1\\
	&= \begin{pmatrix}i\\ -1\\ 0\end{pmatrix}- \frac{3i}{14}\begin{pmatrix}1\\ 2i\\ 3\end{pmatrix}\\
	&= \begin{pmatrix}i- \frac{3}{14}i\\ -1+ \frac{3}{7}\\ -\frac{1}{14}i\end{pmatrix}\\
	&= \begin{pmatrix}\frac{11}{14}i\\ -\frac{4}{7}\\ -\frac{1}{14}i\end{pmatrix}\\
	&= \begin{pmatrix}11i\\ -4\\ -i\end{pmatrix}
\end{align*}
\end{enumerate}
\end{souanswer}

\begin{tcolorbox} %%R_4_4rd_3
	$\fbox{3}$ $C^2$級関数$f(x, y)$に対して、$g(s, t)= f(e^s \cos{t}, e^s \sin{t})$と定める。以下$f_x= \frac{\partial f}{\partial x}(e^s \cos{t}, e^s \sin{t})$とし、$f_y, f_{xx}, f{xy}, f_{yy}$についても同様とする。
\begin{enumerate}
	\item 1階偏導関数$\frac{\partial g}{\partial s},\ \frac{\partial g}{\partial t}$を$f_x, f_y, s, t$を用いて表せ。
	\item 2階偏導関数$\frac{\partial^2 g}{\partial s^2}$を$f_x, f_y, f_{xx}, f_{xy}, f_{yx}, s, t$を用いて表せ。
	\item $\frac{\partial^2 g}{\partial s^2}+ \frac{\partial^2 g}{\partial t^2}$を$f_x, f_y, f_{xx}, f_{xy}, f_{yy}, s, t$を用いて簡単に表せ。
	\item $x, y$の多項式でない$f(x, y)$で$\frac{\partial^2 f}{\partial x^2}+ \frac{\partial^2 f}{\partial y^2}$が零関数になるような例を1つ挙げよ。
\end{enumerate}
\end{tcolorbox}

\begin{souanswer} %R4_3rd_3_ans
	$x= e^s\cos{t}, y= e^s\sin{t}$とおく。各変数の偏導関数は以下の通り
\begin{itemize}
	\item $\frac{\partial x}{\partial s}= e^s\cos{t}= x$
	\item $\frac{\partial x}{\partial t}= -e^s\sin{t}= -y$
	\item $\frac{\partial y}{\partial s}= e^s\sin{t}= y$
	\item $\frac{\partial y}{\partial t}= e^s\cos{t}= x$
\end{itemize}
\begin{enumerate}
	\item 連鎖律を用いる。
\begin{equation*}
	\frac{\partial g}{\partial s}= \frac{\partial f}{\partial x}\frac{\partial x}{\partial s}+ \frac{\partial f}{\partial y}\frac{\partial y}{\partial s}= f_x(e^s\cos{t})+ f_y(e^s\sin{t})
\end{equation*}
\begin{equation*}
	\frac{\partial g}{\partial t}= \frac{\partial f}{\partial x}\frac{\partial x}{\partial t}+ \frac{\partial f}{\partial y}\frac{\partial y}{\partial t}= f_x(-e^s\sin{t})+ f_y(e^s\cos{t})
\end{equation*}
	\item さらに$s$で微分する。積の微分法と連鎖律を用いる
\begin{equation*}
	\frac{\partial^2 g}{\partial s^2}= \frac{\partial}{\partial s}(e^s(f_x\cos{t}+ f_y\sin{t}))= e^s(f_x\cos{t}+ f_y\sin{t})+ e^s\left(\frac{\partial f_x}{\partial s}\cos{t}+ \frac{\partial f_y}{\partial s}\sin{t}\right)
\end{equation*}
	ここで$\frac{\partial f_x}{\partial s}= f_{xx}\frac{\partial x}{\partial s}+ f_{xy}\frac{\partial y}{\partial s}= f_{xx}(e^s\cos{t})+ f_{xy}(e^s\sin{t})$である。
\begin{equation*}
	\frac{\partial^2 g}{\partial s^2}= e^s(f_x\cos{t} f_y\sin{t})+ e^{2s}(f_{xx}\cos^2{t}+ 2f_{xy}\sin{t}\cos{t}+ f_{yy}\sin^2{t})
\end{equation*}

	\item 同様に$\frac{\partial^2 g}{\partial t^2}$を計算し、足し合わせると
\begin{equation*}
	\frac{\partial^2 g}{\partial s^2}+ \frac{\partial^2 g}{\partial t^2}= e^{2s}(f_{xx}+ f_{yy})
\end{equation*}

	\item $\frac{\partial^2 g}{\partial s^2}+ \frac{\partial^2 g}{\partial t^2}= e^{2s}(f_{xx}+ f_{yy})$を満たす関数を調和関数という。
\begin{itemize}
	\item 指数関数と三角関数：$f(x, y)= e^x\sin{y},\ f(x, y)= e^x\cos{y}$
	\item 対数関数：$f(x, y)= \log{(x^2+ y^2)}$　ただし原点を除く
\end{itemize}

\end{enumerate}
\end{souanswer}


\begin{tcolorbox} %%R4_3rd_4
	$\fbox{4}$ $xy$平面上の関数$f(x, y)$を
\begin{equation*}
	f(x, y)= 
\begin{cases}
\begin{matrix}
	\frac{x^2+ xy+ y^2}{\sqrt{x^2+ y^2}}& (x, y)\ne (0 ,0)\\
	0& (x, y)= (0, 0)
\end{matrix}
\end{cases}
\end{equation*}
	で定義する。
\begin{enumerate}
	\item $f(x, y)$は原点$(0, 0)$で連続であることを示せ。
	\item $R$を正の実数とし、$D_R= \{(x, y)\in \R^2; x^2+ y^2\le R^2\}$とおくとき、次の重積分を求めよ。
\begin{equation*}
	\iint_{D_R} f(x, y) dxdy
\end{equation*}
\end{enumerate}
\end{tcolorbox}

\begin{souanswer} %R4_3rd_4_ans
	$\fbox{4}$
\begin{enumerate}
	\item 関数$f(x, y)$が原点で連続であることは、
\begin{equation*}
	\lim_{(x, y)\to (0, 0)} f(x, y)= f(0, 0)= 0
\end{equation*}
	が成り立つことである。極座標$x= r\cos{\theta}, y= r\sin{\theta}$を導入すると$(x, y)\to (0, 0)$は$r\to 0$と同値。$f(x, y)$と書き換えると、
\begin{align*}
	f(x, y)
	&= f(r\cos{\theta}, r\sin{\theta})\\
	&= \frac{r^2\sin^2{\theta}+ r\sin{\theta}\cos{\theta}+ r^2\cos^2{\theta}}{\sqrt{r^2(\cos^2{\theta}+ \sin^2{\theta})}}\\
	&= \frac{r^2(1+ \sin{\theta}\cos{\theta})}{r}\\
	&= r(1+ \sin{\theta}\cos{\theta})
\end{align*}
	ここで$\sin{\theta}\cos{\theta}= \frac{1}{2}\sin{2\theta}$より
\begin{equation*}
	\begin{vmatrix}f(x, y)\end{vmatrix}= \begin{vmatrix}r\left(1+ \frac{1}{2}\sin{2\theta}\right)\end{vmatrix}\le r\left(1+ \frac{1}{2}\right)= \frac{3}{2}r
\end{equation*}
	$r= \sqrt{x^2+ y^2}\to 0$のとき$\frac{3}{2}r\to 0$となるので、はさみうちの原理より
\begin{equation*}
	\lim_{(x, y)\to (0, 0)} f(x, y)= 0
\end{equation*}
	よって原点で連続。

	\item 積分範囲は以下の通り
\begin{itemize}
	\item $0\le r\le R$
	\item $0\le \theta\le 2\pi$
\end{itemize}
\begin{align*}
	\iint_{D_R} f(x, y) dxdy
	&= \int_{\theta= 0}^{2\pi}\int_{r= 0}^R r\left(1+ \frac{1}{2}\sin{2\theta}\right)\cdot rdrd\theta\\
	&= \int_{\theta= 0}^{2\pi} \left(1+ \frac{1}{2}\sin{2\theta}\right)d\theta\cdot \int_{r= 0}^R r^2 dr\\
	&= \left[\theta- \frac{1}{4}\cos{2\theta}\right]_{\theta= 0}^{2\pi}\cdot \left[\frac{1}{3}r^3\right]_{r= 0}^ R\\
	&= \frac{2}{3}\pi R^3
\end{align*}
\end{enumerate}
\end{souanswer}



\chapter{令和5年度}
\section{一般第1期}
	志願者不在の為実施なし

\section{一般第2期}
\begin{tcolorbox} %%R5_2nd_1
	$\fbox{1}$ 次の問いに答えなさい。
\begin{enumerate}
	\item ベクトル$\bm{x}= \begin{pmatrix}7\\ 5\\ 3\end{pmatrix}$をベクトル$\bm{a}_1= \begin{pmatrix}1\\ 2\\ -2\end{pmatrix}$と$\bm{a}_2= \begin{pmatrix}4\\ -1\\ 3\end{pmatrix}$の線形結合で表せ。表せない場合はその理由を説明せよ。
	\item 実数成分行列$A= \begin{pmatrix}x& 1& 1\\ 1& x& 1\\ 1& 1& x\end{pmatrix}$の階数を$x$の値によって場合分けして求めなさい。
	\item $\R^3$を実数体$\R$の3次元ベクトル空間とし、写像$f\colon \R^3\to \R^3$を$f\left(\begin{pmatrix}x\\ y\\ z\end{pmatrix}\right)= \begin{pmatrix}{x+ 2y- z}\\ {y+ z}\\ {x+ y- 2z}\end{pmatrix}$と定義する。
\begin{enumerate_roman}
	\item $f$が線形写像であることを示せ。
	\item $f$の核の次元を求めよ。
\end{enumerate_roman}
\end{enumerate}
\end{tcolorbox}

\begin{souanswer} %%R5_2nd_1_ans
	$\fbox{1}$
\begin{enumerate}
	\item 線型結合で表せるならば$\bm{x}= c_1\bm{a}_1+ c_2\bm{a}_2$となる実数係数$c_1, c_2$を見つける。
\begin{equation*}
	c_1\begin{pmatrix}1\\ 2\\ -2\end{pmatrix}+ c_2\begin{pmatrix}4\\ -1\\ 3\end{pmatrix}= \begin{pmatrix}7\\ 5\\ 3\end{pmatrix}
	\iff
	\begin{cases}
		c_1+ 4c_2= 7\\
		2c_1- c_2= 5\\
		-2c_1+ 3c_2= 3
	\end{cases}
\end{equation*}
	を解く。これを満たす$c_1, c_2$は存在しない。よって$\bm{x}$を$\bm{a}_1, \bm{a}_2$の線型結合で表せない(互いに線型独立)。
	\item 行列$A$の階数を調べるために行列式$|A|$を計算し、正則性を調べる。
\begin{equation*}
	|A|= \begin{vmatrix}x& 1& 1\\ 1& x& 1\\ 1& 1& x\end{vmatrix}= (x+2)(x-1)^2= 0\iff x= -2, 1
\end{equation*}
	よって以下の\ref{pattern:xn1andxn-2}, \ref{pattern:x=1} , \ref{pattern:x=1} の3パターンに分けられる。
\begin{enumerate_roman}
	\item\label{pattern:xn1andxn-2} $x\ne 1$かつ$x\ne -2$のとき：行列式が0ではないため、行列$A$は正則。したがって$\rank{A}= 3$
	\item\label{pattern:x=1} $x= 1$のとき：$A= \begin{pmatrix}1& 1& 1\\ 1& 1& 1\\ 1& 1& 1\end{pmatrix}$となり、独立な行は1つだけなので$\rank{A}= 1$
	\item\label{pattern:x=2} $x= -2$のとき：$A= \begin{pmatrix}-2& 1& 1\\ 1& -2& 1\\ 1& 1& -2\end{pmatrix}$となり、第1列と第2列は独立だが、それぞれの行の和は0になるので、3つの行ベクトルは従属。$\rank{A}= 2$
\end{enumerate_roman}
\begin{enumerate}
	\item 写像$f$が線型写像であることを示すためには任意のベクトル$\bm{u}, \bm{v}\in \R^3$とスカラー$k\in \R$に対して以下の条件\ref{law:加法性}, \ref{law:斉次性}
\begin{enumerate_roman}
	\item\label{law:加法性} $f(\bm{u}+ \bm{v})= f(\bm{u})+ f(\bm{v})$
	\item\label{law:斉次性} $f(k\bm{u})= kf(\bm{u})$
\end{enumerate_roman}
	この写像は行列$M= \begin{pmatrix}1& 2& -1\\ 0& 1& 1\\ 1& 1& -2\end{pmatrix}$を用いて$f(\bm{v})= M\bm{v}$と表現できる。行列とベクトルの積の性質より
\begin{enumerate}
	\item $M(\bm{u}+ \bm{v})= M\bm{u}+ M\bm{v}$
	\item $M(k\bm{u})= kM\bm{u}$
\end{enumerate}
	が成り立つため、$f$は線型写像。
\end{enumerate}
	\item 核の次元を求めるために、表現行列$M$の階数を調べる。
\begin{equation*}
	M= \begin{pmatrix}1& 2& -1\\ 0& 1& 1\\ 1& 1& -2\end{pmatrix}
\stackrel{R_3- R_1}{\longrightarrow} \begin{pmatrix}1& 2& -1\\ 0& 1& 1\\ 0& -1& -1\end{pmatrix}\stackrel{R_3+ R_2}{\longrightarrow}
\begin{pmatrix}1& 2& -1\\ 0& 1& 1\\ 0& 0& 0\end{pmatrix}
\end{equation*}
	階段行列の形から$\rank{M}= 2$であることがわかる。次元定理より
\begin{equation*}
	\dim{\ker{f}}= 3- 2= 1
\end{equation*}
	核の次元は1
\end{enumerate}
\end{souanswer}

\begin{tcolorbox} %%R5_2nd_2
	$\fbox{2}$ 次の問いに答えなさい。
\begin{enumerate}
	\item $A$を$n$次複素正方行列とし、$P$を$n$次複素正則行列とする。このとき複素数$\lambda$が$A$の固有値ならば、$\lambda$は$P^{-1}AP$の固有値でもあることを示せ。
	\item $V$を複素数上の有限次元線型空間とし、$W$をその部分空間とする。また$V$には内積$\innerproduct{\bullet}{\bullet}$が備わっているとする。
\begin{enumerate_roman}
	\item $W$への射影子$P$がどのように定義されるか、用いる定理を明記して説明せよ。
	\item $P$を$W$への射影子としたとき、任意の$x, y\in V$に対して、${Px}{y}= {x}{Py}$が成り立つことを示せ。
\end{enumerate_roman}
\end{enumerate}
\end{tcolorbox}

\begin{souanswer} %%R5_2nd_2
	$\fbox{2}$
\begin{enumerate}
	\item 行列$B= P^{-1}AP$の固有多項式を計算すると
\begin{equation*}
	\begin{vmatrix}B- \lambda E\end{vmatrix}= \begin{vmatrix}P^{-1}AP- \lambda P^{-1}EP\end{vmatrix}= \begin{vmatrix}P^{-1}(A- \lambda E)P\end{vmatrix}
\end{equation*}
	ここで行列式の性質($|XY|= |X||Y|$)より
\begin{align*}
	\begin{vmatrix}B- \lambda E\end{vmatrix}
	&= \begin{vmatrix}P^{-1}AP- \lambda E\end{vmatrix}\\
	&= \begin{vmatrix}P^{-1}\end{vmatrix}\begin{vmatrix}A\end{vmatrix}\begin{vmatrix}P\end{vmatrix}\\
	&= \frac{1}{\begin{vmatrix}P\end{vmatrix}}\cdot\begin{vmatrix}A- \lambda E\end{vmatrix}\cdot \begin{vmatrix}P\end{vmatrix}\\
	&= \begin{bmatrix}A- \lambda E\end{bmatrix}
\end{align*}
	$\lambda$は$A$の固有値だから$\begin{vmatrix}A- \lambda E\end{vmatrix}= 0$。したがって$\begin{vmatrix}B- \lambda E\end{vmatrix}= 0$となるので、$\lambda$は$B$の固有値でもある。
	\item 
\begin{enumerate_roman}
	\item $V$を有限次元内積空間、$V$をその部分空間とする。このとき直交補空間の存在定理
\begin{tcolorbox}[title= 直交補空間の存在定理]
	$V$の任意の真要素$v$は$w\in W$と$w^\perp\in W^\perp$($W$の直交補空間)を用いて、$v= w+ w^\perp$と一意に分解できる。すなわち$V= W\oplus W^\perp$が成り立つ。
\end{tcolorbox}
	この定理に基づき、$v\in V$に対して$w\in W$を対応させる写像$P\colon V\to V$を$W$への直交射影子と定義する。

	\item 直和分解$V= W\oplus W^\perp$により任意の$x, y\in V$を次のように分割する。
\begin{itemize}
	\item $x= x_w+ x_{w^\perp}\ \ (w\in W, w^\perp\in W^\perp)$
	\item $y= y_w+ y_{w^\perp}\ \ (w\in W, w^\perp\in W^\perp)$
\end{itemize}
	射影の定義より$Px= x_w, Py= y_w$である。内積の性質(線型性と直交性)より
\begin{align*}
	(左辺)
	&= \innerproduct{x}{Py}\\	
	&= \innerproduct{x_w}{y_w+ y_{w^\perp}}\\
	&= \innerproduct{x_w}{y_w}+ \innerproduct{x_w}{y_{w^\perp}}\\
	&= \innerproduct{x_w}{y_w}\ \ (\because x_w\perp y_{w^\perp})\\
	(右辺)
	&= \innerproduct{Px}{y}\\	
	&= \innerproduct{x_w+ x_{w^\perp}}{y_w}\\
	&= \innerproduct{x_w}{y_w}+ \innerproduct{x_{w^\perp}}{y_w}\\
	&= \innerproduct{x_w}{y_w}\ \ (\because x_{w^\perp}\perp y_w)\\
\end{align*}
	よって$\innerproduct{Px}{y}= \innerproduct{x}{Py}$が成り立つ。
\end{enumerate_roman}
\end{enumerate}
\end{souanswer}

\begin{tcolorbox} %%R5_2nd_3
	$\fbox{3}$ 自然数$n$に対して、閉区間$[-1, 1]$上の関数$f_n$を
\begin{equation*}
	f_n(x)=
\begin{cases}
\begin{matrix}
	n& -\frac{1}{2n}\le x\le \frac{1}{2n}\\
	0& x< -\frac{1}{2n}あるいは x> \frac{1}{2n}
\end{matrix}
\end{cases}
\end{equation*}
	で定義する。
\begin{enumerate}
	\item $y= f_1(x), y= f_2(x), y= f_3(x)$のグラフを図示せよ。
	\item 自然数$n$に対して$\int_{-1}^1 f_n(x) dx$を求めよ。
	\item $n$が自然数で$g\colon \R\to \R$が連続的微分可能な関数のとき$\int_{-1}^1 f_n(x)g'(x) dx$を求めよ。
	\item $g\colon \R\to \R$が連続的微分可能な関数のとき、$\lim_{n\to \infty}\int _{-1}^1 f_n(x)g'(x) dx= g'(0)$を示せ。
\end{enumerate}
\end{tcolorbox}

\begin{souanswer} %%R5_2nd_3_ans
	$\fbox{3}$
\begin{enumerate}
	\item $y= f_1(x), y= f_2(x), y= f_3(x)$それぞれのグラフを描く。
%\begin{figure}[h]
\begin{tikzpicture}%f_1のグラフ
	\draw[->, >= stealth, semithick] (-2, 0)-- (2, 0)node[above]{$x$}; %x軸
	\draw[->, >= stealth, semithick] (0, -1)-- (0, 3.5)node[right]{$y$}; %y軸
	\draw (0, 0) node[below right]{O}; %原点
	\draw (-1/2, 0)node[below left]{$-\frac{1}{2}$}; %点(-1/2, 0)
	\draw (1/2, 0)node[below right]{$\frac{1}{2}$}; %点(0, 1/2)
	\draw[very thick, domain= -1/2: 1/2] plot(\x, 1)node[above right]{$y= f_1(x)$}; %f_1のグラフ
	\draw[dashed] (-1/2, 0)-- (-1/2, 1);
	\draw[dashed] (1/2, 0)-- (1/2, 1);
\end{tikzpicture}
%	\caption{$y= f_1(x)$のグラフ}
%\end{figure}

%\begin{figure}[h]
\begin{tikzpicture}%f_2のグラフ
	\draw[->, >= stealth, semithick] (-2, 0)-- (2, 0)node[above]{$x$}; %x軸
	\draw[->, >= stealth, semithick] (0, -1)-- (0, 3.5)node[right]{$y$}; %y軸
	\draw (0, 0) node[below right]{O}; %原点
	\draw (-1/4, 0)node[below left]{$-\frac{1}{4}$}; %点(-1/4, 0)
	\draw (1/4, 0)node[below right]{$\frac{1}{4}$}; %点(0, 1/4)
	\draw[very thick, domain= -1/4: 1/4] plot(\x, 2)node[above right]{$y= f_2(x)$}; %f_2のグラフ
	\draw[dashed] (-1/4, 0)-- (-1/4, 2);
	\draw[dashed] (1/4, 0)-- (1/4, 2);
\end{tikzpicture}
%	\caption{$y= f_2(x)$のグラフ}
%\end{figure}

%\begin{figure}[h]
\begin{tikzpicture}%f_3のグラフ
	\draw[->, >= stealth, semithick] (-2, 0)-- (2, 0)node[above]{$x$}; %x軸
	\draw[->, >= stealth, semithick] (0, -1)-- (0, 3.5)node[right]{$y$}; %y軸
	\draw (0, 0) node[below right]{O}; %原点
	\draw (-1/6, 0)node[below left]{$-\frac{1}{6}$}; %点(-1/9, 0)
	\draw (1/6, 0)node[below right]{$\frac{1}{6}$}; %点(0, 1/9)
	\draw[very thick, domain= -1/6: 1/6] plot(\x, 3)node[above right]{$y= f_3(x)$}; %f_2のグラフ
	\draw[dashed] (-1/6, 0)-- (-1/6, 3);
	\draw[dashed] (1/6, 0)-- (1/6, 3);
\end{tikzpicture}
%	\caption{$y= f_3(x)$のグラフ}
%\end{figure}
	\item 定義域を関数値の0でない区間に限定して積分する。
\begin{align*}
	\int_{-1}^1 f_n(x) dx
	&= \int_{-\frac{1}{2n}}^{\frac{1}{2}} n dx\\
	&= n[x]_{-\frac{1}{2n}}^{\frac{1}{2}}\\
	&= n\left(\frac{1}{2n}- \left(-\frac{1}{2n}\right)\right)\\
	&= n\cdot \frac{1}{n}\\
	&= 1
\end{align*}
	\item\label{val:g'} 同様に積分範囲を限定して計算する。
\begin{align*}
	\int_{-1}^1 f_n(x)g'(x) dx
	&= \int_{-\frac{1}{2n}}^\frac{1}{2n} n\cdot g'(x) dx\\
	&= n\int_{-\frac{1}{2n}}^\frac{1}{2n} g'(x) dx\\
	&= n[g'(x)]_{-\frac{1}{2n}}^\frac{1}{2n}\\
	&= n\left(g\left(\frac{1}{2n}\right)- g\left(-\frac{1}{2n}\right)\right)
\end{align*}
	\item (\ref{val:g'})より
\begin{equation*}
	\lim_{n\to \infty}\int_{-1}^1 f_n(x)g'(x) dx= \lim_{n\to \infty} n\left(g\left(\frac{1}{2n}\right)- g\left(-\frac{1}{2n}\right)\right)
\end{equation*}
	ここで$h= \frac{1}{2n}$とおくと、$n\to \infty$のとき$h\to 0$。また$n= \frac{1}{2h}$なので
\begin{equation*}
	\lim_{h\to 0} \frac{1}{2h}(g(h)- g(-h))= \lim_{h\to 0}\frac{g(h)- g(-h)}{2h}
\end{equation*}
	ここで$-g(0)+ g(0)$を挿入し、整理すると
\begin{equation*}
	\lim_{h\to 0} \frac{1}{2}\left(\frac{g(h)- g(0)}{h}+ \frac{g(0)- g(-h)}{h}\right)= \frac{1}{2}(g'(0)+ g'(0))= g'(0)
\end{equation*}
	$g$は連続的微分可能であるため、$g'(0)$が存在し、この極限は確定する。
\end{enumerate}
\end{souanswer}

\begin{tcolorbox} %%R5_2nd_4
	$\fbox{4}$ 正の実数$\alpha, t$と実数$x, y, z$に対して$\rho_\alpha(x, y, z, t)= \frac{1}{(4\pi t)^\alpha}e^{-\frac{x^2+ y^2+ z^2}{4t}}$とおく。
\begin{enumerate}
	\item $\frac{\partial \rho_\alpha}{\partial x}$を求めよ。
	\item $\frac{\partial^2 \rho_\alpha}{\partial^2 x}$を求めよ。
	\item $\frac{\partial \rho_\alpha}{\partial t}$を求めよ。
	\item $\frac{\partial \rho_\alpha}{\partial t}= \frac{\partial^2 \rho_\alpha}{\partial x^2}+ \frac{\partial^2 \rho_\alpha}{\partial y^2}+ \frac{\partial^2 \rho_\alpha}{\partial z^2}$となるための$\alpha$の条件を求めよ。
\end{enumerate}
\end{tcolorbox}
\begin{souanswer} %%R5_2nd_4_ans
	$\fbox{4}$
\begin{enumerate}
	\item 
\begin{align*}
	\frac{\partial \rho_\alpha}{\partial x}
	&= \frac{1}{(4\pi t)^\alpha}e^{-\frac{x^2+ y^2+ z^2}{4t}}\cdot \frac{\partial}{\partial x}\left(-\frac{x^2+ y^2+ z^2}{4t}\right)\\
	&= \rho_\alpha\cdot\left(-\frac{2x}{4t}\right)\\
	&= -\frac{x}{2t}\rho_\alpha
\end{align*}

	\item 積の微分より
\begin{align*}
	\frac{\partial^2 \rho_\alpha}{\partial x^2}
	&= \frac{\partial}{\partial x}\left(-\frac{x}{2t}\rho_\alpha\right)\\
	&= -\frac{1}{2t}\rho_\alpha- \frac{x}{2t}\cdot \frac{\partial}{\partial x}\rho_\alpha\\
	&= -\frac{1}{2t}\rho_\alpha+ \frac{x^2}{4t^2}\rho_\alpha\\
	&= \left(\frac{x^2}{4t^2}- \frac{1}{2t}\right)\rho_\alpha
\end{align*}

	\item $t$については積の微分を用いる
\begin{equation*}
	\frac{\partial \rho_\alpha}{\partial t}
	= \frac{\partial}{\partial t}((4\pi t)^{-\alpha})\cdot e^{-\frac{x^2+ y^2+ z^2}{4t}}+ (4\pi t)^{-\alpha}\cdot \frac{\partial}{\partial t}\left(e^{-\frac{x^2+ y^2+ z^2}{4t}}\right)
\end{equation*}
\begin{align*}
	(第一項)
	&= -\alpha(4\pi)^{-\alpha}t^{-\alpha- 1}e^{-\frac{x^2+ y^2+ z^2}{4t}}\\
	&= -\frac{\alpha}{t}\rho_\alpha\\
	(第二項)
	&= \rho_\alpha\cdot \frac{\partial}{\partial t}\left({-\frac{x^2+ y^2+ z^2}{4t}}\right)\\
	&= \rho_\alpha\left({-\frac{x^2+ y^2+ z^2}{4t}}\right)
\end{align*}
	よって
\begin{equation*}
	\frac{\partial \rho_\alpha}{\partial t}= \left(\frac{x^2+ y^2+ z^2}{4t^2}- \frac{\alpha}{t}\right)\rho_\alpha
\end{equation*}

	\item 与えられた方程式は以下(\eqref{eq:熱方程式}とする。)
\begin{equation} \tag{H}\label{eq:熱方程式}
	\frac{\partial \rho_\alpha}{\partial t}= \Delta\rho_\alpha
\end{equation}
	\eqref{eq:熱方程式}の右辺のLaplacianを計算する。
\begin{equation*}
	\frac{\partial^2 \rho_\alpha}{\partial y^2}= \left(\frac{y^2}{4t^2}- \frac{1}{2t}\right)\rho_\alpha,\ 
	\frac{\partial^2 \rho_\alpha}{\partial z^2}= \left(\frac{z^2}{4t^2}- \frac{1}{2t}\right)\rho_\alpha
\end{equation*}
	これらを足し合わせて
\begin{equation*}
	\Delta\rho_\alpha= \left(\frac{x^2+ y^2+ z^2}{4t}- \frac{3}{2t}\right)\rho_\alpha
\end{equation*}
	求めた左辺と比較して
\begin{equation*}
	\left(\frac{x^2+ y^2+ z^2}{4t^2}- \frac{\alpha}{t}\right)\rho_\alpha= \left(\frac{x^2+ y^2+ z^2}{4t}- \frac{3}{2t}\right)\rho_\alpha
\end{equation*}
	係数比較して$-\frac{\alpha}{t}= -\frac{1}{2t}\iff \alpha= \frac{3}{2}$

\end{enumerate}
\end{souanswer}

\section{一般第3期}
	志願者不在の為実施なし



\chapter{令和6年度}
\section{一般第1期}
\begin{tcolorbox} %%R6_1st_1
	$\fbox{1}$ 次の問いに答えよ。
\begin{enumerate}
	\item $V$を実数体$\R$上の線形空間とする。$V$の空でない部分集合$W$が$V$の部分空間であることの定義を書け。
	\item 実数体$\R$上の$n$次元ベクトル空間$\R^n$に対して、次の集合は部分空間となるかどうか理由を付けて判定せよ。ただし$\bm{v}\in \R^n$の第$i$成分を$v_i$として表す。
\begin{enumerate_roman}
	\item $W_1=\{\bm{v}\in \R^n; v_1+ v_2+ \cdots+ v_n= 0\}$
	\item $W_2= \{\bm{v}\in \R^n; v_1\times v_2\times \cdots\times v_n= 0\}$
\end{enumerate_roman}

	\item 定義域を$\R$とする実数値関数$f\colon \R\to \R$の集合$V$とする。$V$は実数体$\R$の線形空間である(これは示さなくてよい)。次の集合は$V$の部分空間になるかどうか理由をつけて判定せよ。
\begin{enumerate_roman}
	\item $U_1=\{f\in V; fは有界\}$
	\item $U_2= \{f\in V; fは連続でない\}$
\end{enumerate_roman}
\end{enumerate}
\end{tcolorbox}

\begin{souanswer} %%R6_1st_1_ans
	$\fbox{1}$
\begin{enumerate}
	\item $V$の任意の元$x\in W$に対して
\begin{equation*}
	W= \{x\in V; \innerproduct{x}{a},\ \forall a\in W\}
\end{equation*}
	\item 
\begin{enumerate_roman}
	\item 
	\item 
\end{enumerate_roman}
	\item 
\begin{enumerate_roman}
	\item 
	\item 
\end{enumerate_roman}
\end{enumerate}
\end{souanswer}




\begin{tcolorbox} %%R6_1st_2
	$\fbox{2}$ $A= \begin{pmatrix}1& i& i\\ i& i& 1\end{pmatrix}$とし、$X$を複素行列とする。
\begin{enumerate}
	\item $A$の随伴行列$A^{*}$を求めよ。
	\item $A^{*}A$を計算せよ。
	\item エルミート行列の定義を述べよ。また$X^{*}X$がエルミート行列であることを示せ。
\end{enumerate}
\end{tcolorbox}

\begin{souanswer} %%R6_1st_2_ans
\begin{enumerate}
	\item 随伴行列$A^{*}$は
\begin{equation*}
	A^{*}= 
\begin{pmatrix}
	\bar{1}& \bar{i}\\
	\bar{i}& \bar{i}\\
	\bar{i}& \bar{1}
\end{pmatrix}
	=
\begin{pmatrix}
	1& -i\\
	-i& -i\\
	-i& 1
\end{pmatrix}
\end{equation*}
	\item
\begin{align*}
	A^{*}A
	&=
\begin{pmatrix}
	1& -i\\
	-i& -i\\
	-i& 1
\end{pmatrix}
\begin{pmatrix}
	1& i& i\\
	i& i& 1
\end{pmatrix}\\
	&= 
\begin{pmatrix}
	2& 1+ i& 0\\
	1- i& 2& 1- i\\
	0& 1+ i& 2
\end{pmatrix}
\end{align*}
	\item エルミート行列の定義は以下
\begin{tcolorbox}
	随伴行列と元の行列が等しい正方行列。すなわち正方行列$A$に対して$A^*= \bar{A}^\mathrm{T}= A$となる行列。
\end{tcolorbox}
	任意の複素行列$X$に対して$H= X^*X$とおく。その随伴行列を$H^*$を計算する。随伴行列の性質$(AB)^*= B^*A^*$および$(X^*)^*= X$より
\begin{equation*}
	(X^*X)^*= X^*(X^*)^*= X^*X
\end{equation*}
	$(X^*X)^*= X^*X$が成り立つため$X^*X$はエルミート行列

\end{enumerate}
\end{souanswer}

\begin{tcolorbox} %%R6_1st_3
	$\fbox{3}$ 
$f(x, y)= 
\begin{cases}
\begin{matrix}
	\frac{\sin{(xy(2x^2- 3y^2))}}{x^2+ y^2}& (x, y)\ne (0, 0)\\
	0& (x, y)= (0, 0)
\end{matrix}
\end{cases}
$
	に対して次の問いに答えよ。
\begin{enumerate}
	\item $\frac{\partial f}{\partial x}(x, y),\ \frac{\partial}{\partial y}\frac{\partial f}{\partial x}(0, 0)$を求めよ。
	\item $\frac{\partial f}{\partial y}(x, y),\ \frac{\partial}{\partial y}\frac{\partial f}{\partial x}(0, 0)$を求めよ。
	\item $f(x, y)$は$(0, 0)$において全微分可能かどうか理由をつけて判定せよ。
\end{enumerate}
\end{tcolorbox}

\begin{souanswer} %%R6_1st_3_ans
	$\fbox{3}$
\begin{enumerate}
	\item 商の微分を用いて計算する。$u= xy(2x^2- 3y^2)= 2x^3y- 3xy^3$とおくと、$\partial u/\partial x= 6x^2y- 3y^3$。
\begin{equation*}
	\frac{\partial f}{\partial x}= \frac{(6x^2y- 3y^2)(x^2+ y^2)\cos{u}- 2x\sin{u}}{(x^2+ y^2)^2}
\end{equation*}
	偏微分の定義より
\begin{equation*}
	\frac{\partial f}{\partial x}(0, y)= \lim_{h\to 0} \frac{f(h, 0)- f(0, y)}{h}= \lim_{h\to 0} \frac{\frac{\sin{(hy(2h^2- 3y^2))}}{h^2+ y^2}- 0}{h}
\end{equation*}
	ここで$\sin{\theta}\to 0(\theta\to 0)$より
\begin{equation*}
	\frac{\partial f}{\partial x}(0, y)= \lim_{h\to 0} \frac{hy(2h^2- 3y^2)}{h(h^2+ y^2)}= \frac{-3y^3}{y^2}= -3y
\end{equation*}
	これを用いて$(0, 0)$における$y$に関する偏微分を求める。
\begin{align*}
	\frac{\partial^2 f}{\partial y\partial x}(0, 0)
	&= \lim_{k\to 0} \frac{\frac{\partial f}{\partial x}(0, k)- \frac{\partial f}{\partial x}(0, 0)}{k}\\
	&=\lim_{k\to 0} \frac{-3k- 0}{k}= -3
\end{align*}

	\item 同様に$\frac{\partial u}{\partial y}= 2x^3-9xy^2$なので
\begin{equation*}
	\frac{\partial f}{\partial y}= \frac{(2x^3- 9xy^2)(x^2+ y^2)\cos{u}- 2y\sin{u}}{(x^2+ y^2)^2}
\end{equation*}
	$\frac{\partial f}{\partial y}(x, 0)$を求める。
\begin{align*}
	\frac{\partial f}{\partial y}(x, 0)
	&= \lim_{k\to 0} \frac{f(x, k)- f(x, 0)}{k}\\
	&= \lim_{k\to 0} \frac{\frac{\sin(xk(2x^2- 3k^2))}{x^2+ k^2}- 0}{k}\\
	&= \frac{2x^3}{x^2}\\
	&= 2x
\end{align*}
	これを用いて
\begin{equation*}
	\frac{\partial^2 f}{\partial x\partial y}(0, 0)= \lim_{h\to 0} \frac{\frac{\partial f}{\partial y}(h, 0)- \frac{\partial f}{\partial y}(0, 0)}{h}= \lim_{h\to 0} \frac{2h- 0}{h}= 2
\end{equation*}

	\item (0, 0)において全微分可能であるためには以下の極限が0に収束する必要がある。
\begin{equation*}
	\lim_{(h, k)\to (0, 0)} \frac{f(h, k)- f(0, 0)- \left(h\frac{\partial f}{\partial x}(0, 0)+ k\frac{\partial f}{\partial y}(0, 0)\right)}{\sqrt{h^2+ k^2}}= 0
\end{equation*}
	$\frac{\partial f}{\partial x}(0, 0), \frac{\partial f}{\partial y}(0, 0)$はいずれも0。判定式は
\begin{equation*}
	\epsilon= \frac{\frac{\sin(hk(2h^2- 3k^2))}{h^2+ k^2}}{\sqrt{h^2+ k^2}}= \frac{\sin{(hk(2h^2- 3k^2))}}{(h^2+ k^2)^{3/2}}
\end{equation*}
	ここで極座標$h= r\cos{\theta}, k= r\sin{\theta}$を用いると分子の引数は$O(r^4)$となる。
\begin{equation*}
	\epsilon\sim \frac{r^2\cos{\theta}\sin{\theta}(2r^2\cos^2{\theta}- 3r^2\sin^2{\theta})}{r^3}= r\cos{\theta}\sin{\theta}(2\cos^2{\theta}- 3\sin^2{\theta})
\end{equation*}
	$r\to 0$のときこの値は$\theta$によらず0に収束する。
\end{enumerate}
\end{souanswer}


\begin{tcolorbox} %%R6_1st_4
	$\fbox{4}$ $a, b, c$を正の定数とし、3次元実空間の原点を$O$とする。$x_0, y_0, z_0$を0でない実数とし、曲面$\frac{x^2}{a^2}+ \frac{y^2}{b^2}+ \frac{z^2}{c^2}= 1$の点$(x_0, y_0, z_0)$における接平面$H$が$x$軸、$y$軸、$z$軸と交わる点をそれぞれX, Y, Zとする。また四面体OXYZの体積を$V$とする。
\begin{enumerate}
	\item 曲面$\frac{x^2}{a^2}+ \frac{y^2}{b^2}+ \frac{z^2}{c^2}= 1$で囲まれる図形の体積を求めよ。
	\item 接平面$H$の式を$(x_0, y_0, z_0)$を用いて書け。
	\item X, Y, Zを$(x_0, y_0, z_0)$を用いて表しなさい。
	\item Vを$(x_0, y_0, z_0)$を用いて表しなさい。
	\item $\frac{x_0^2}{a^2}+ \frac{y_0^2}{b^2}+ \frac{z_0^2}{c^2}= 1$のもとでVの最小値を求めよ。
\end{enumerate}
\end{tcolorbox}

\begin{souanswer} %%R6_1st_4_ans
	$\fbox{4}$
\begin{enumerate}
	\item 体積$V_\mathrm{ell}$は球の体積$V_\mathrm{sp}= \frac{4}{3}\pi r^3$を$x, y, z$方向にそれぞれ$a, b, c$だけ引き延ばしたもの
\begin{equation*}
	V_\mathrm{ell}= \frac{4}{3}\pi abc
\end{equation*}

	\item 点$(x_0, y_0, z_0)$における接平面の方程式は円や球の接線と同様。関数$f(x, y, z)= \frac{x^2}{a^2}+ \frac{y^2}{b^2}+ \frac{z^2}{c^2}-1= 0$の勾配ベクトル$\nabla f$を用いると
\begin{equation*}
	\frac{x_0}{a^2}x+ \frac{y_0}{b^2}y+ \frac{z_0}{c^2}z= 1
\end{equation*}

	\item 接平面$H$と各軸の交点を求める。
\begin{itemize}
	\item $X$点：$y= 0, z= 0$を代入すると$\frac{x_0}{a^2}x= 1\Rightarrow x= \frac{a^2}{x_0}$
	\item $Y$点：$x= 0, z= 0$を代入すると$\frac{y_0}{b^2}y= 1\Rightarrow y= \frac{b^2}{y_0}$
	\item $Z$点：$x= 0, y= 0$を代入すると$\frac{z_0}{c^2}z= 1\Rightarrow x= \frac{c^2}{z_0}$
\end{itemize}

	\item 四面体$\mathrm{OXYZ}$は各軸が直交しているため、底面を$\triangle\mathrm{OXY}$とすると、高さは$\mathrm{OX}$となる。
\begin{align*}
	V
	&= \frac{1}{3}\times (\triangle\mathrm{OXY})\times \mathrm{OZ}\\
	&= \frac{1}{3}\left(\frac{1}{2}\left|\frac{a^2}{x_0}\right|\left|\frac{b^2}{y_0}\right|\right)\left|\frac{c^2}{z_0}\right|\\
	&= \frac{a^2b^2c^2}{6|x_0y_0z_0|}
\end{align*}

	\item 制約条件$\frac{x_0^2}{a^2}+ \frac{y_0^2}{b^2}+ \frac{z_0^2}{c^2}= 1$の元で最小となるためには分母$|x_0y_0z_0|$を最大にする必要がある。相加平均相乗平均より$A= \frac{x_0^2}{a^2}, B= \frac{y_0^2}{b^2}, C= \frac{z_0^2}{c^2}$とおくと$A, B, C> 0$かつ$A+ B+ C= 1$である。
\begin{align*}
	\frac{A+ B+ C}{3}&\ge \sqrt[3]{ABC}\\
	\frac{1}{3}&\ge \sqrt[3]{\frac{x_0^2y_0^2z_0^2}{a^2b^2c^2}}
\end{align*}
	両辺3乗して$\frac{1}{27}\ge \frac{x_0^2y_0^2z_0^2}{a^2b^2c^2}\Rightarrow |x_0^2y_0^2z_0^2|\le \frac{abc}{3\sqrt{3}}$
	$|x_0y_0z_0|$の最大値が$\frac{abc}{3\sqrt{3}}$なので、これを$V$に代入して
\begin{equation*}
	V_\mathrm{min}= \frac{a^2b^2c^2}{6\cdot \frac{abc}{3\sqrt{3}}}= \frac{\sqrt{3}}{2}abc
\end{equation*}
	等号成立は$A= B= C= \frac{1}{3}$。すなわち$x_0= \frac{a}{\sqrt{3}}, y_0= \frac{b}{\sqrt{3}}, z_0= \frac{c}{\sqrt{3}}$のとき

\end{enumerate}
\end{souanswer}

\section{一般第2期}
	志願者不在のため実施無し

\section{一般第3期}
\begin{tcolorbox} %%R6_3rd_1
	$\fbox{1}$ $n$を自然数とし、$A, B, C$を$n\times n$行列とする。また正方行列$M$に対し、$\Tr{(M)}$は$M$の対角成分の総和を表すとする。
\begin{enumerate}
	\item $
	A= 
\begin{pmatrix}
	a_{11}& \cdots& a_{1n}\\
	\vdots& \ddots& \vdots\\
	a_{n1}& \cdots& a_{nn}
\end{pmatrix},\ \ 
	B= 
\begin{pmatrix}
	b_{11}& \cdots& b_{1n}\\
	\vdots& \ddots& \vdots\\
	b_{n1}& \cdots& b_{nn}
\end{pmatrix}$
	とおいたとき、行列$AB$の第$k$対角成分(すなわち$(k, k)$成分)を求めなさい。また行列$BA$の第$k$対角成分を求めなさい。
	\item $\Tr{(AB)}= \Tr{(BA)}$を示せ。
	\item $\Tr{(ABC)}= \Tr{(CAB)}= \Tr{(BCA)}$を示せ。
	\item $P$を$n\times n$正則行列とすると、$\Tr{(P^{-1}AP)}= \Tr{(A)}$が成り立つことを示せ。
	\item $\Tr{(ABC)}= \Tr{(ACB)}$が成り立たない例を挙げよ。
\end{enumerate}
\end{tcolorbox}

\begin{souanswer} %%R6_3rd_1_ans
	$\fbox{1}$ 
\begin{enumerate}
	\item 行列$AB$の対角成分を$(AB)_{kk}$，行列$BA$の対角成分を$(BA)_{kk}$と書く。
\begin{align*}
	(AB)_{kk}
&= \begin{pmatrix}a_{1k}\\ a_{2k}\\ \vdots\\ a_{nk}\end{pmatrix}\begin{pmatrix}b_{k1}& b_{k2}& \cdots& b_{kn}\end{pmatrix}\\
&= a_{1k}b_{k1}+ a_{2k}b_{k2}+ \cdots+ a_{nk}b_{kn}\\
&= \sum_{i= 1}^n a_{ik}b_{ki}
\end{align*}

\begin{align*}
	(BA)_{kk}
&= \begin{pmatrix}b_{1k}\\ b_{2k}\\ \vdots\\ b_{nk}\end{pmatrix}\begin{pmatrix}a_{k1}& a_{k2}& \cdots& a_{kn}\end{pmatrix}\\
&= b_{1k}a_{k1}+ b_{2k}a_{k2}+ \cdots+ b_{nk}a_{kn}\\
&= \sum_{i= 1}^n b_{ik}a_{ki}
\end{align*}

	\item $\Tr{(AB)}= \Tr{(BA)}$を示す。それぞれのトレースを計算すると
\begin{align*}
	\Tr{(AB)}
&= \sum_{k= 1}^n \left(\sum_{j= 1}^ n a_{kj}b_{jk}\right)\\
&= \sum_{k= 1}^n \sum_{j= 1}^ n a_{kj}b_{jk}\\
	\Tr{(BA)}
&= \sum_{k= 1}^n \left(\sum_{j= 1}^ n b_{kj}a_{jk}\right)\\
&= \sum_{k= 1}^n \sum_{j= 1}^n b_{kj}a_{jk}
\end{align*}
	ここで和の順序を変えることができ、$a_{kj}b_{jk}= b_{kj}a_{jk}$であるから$\Tr{(AB)}= \Tr{(BA)}$

	\item 
\begin{equation*}
	\Tr{(AB)}= \sum_{k= 1}^n(AB)_{kk}= \sum_{k= 1}^n \sum_{j= 1}^n a_{kj}b_{jk}
\end{equation*}
	和の順序を入れ替えると
\begin{equation*}
	\Tr{(AB)}= \sum_{j= 1}^n\sum_{k= 1}^n b_{jk}a_{kj}
\end{equation*}
	総和の引数は$BA$の第$j$対角成分$(BA)_{jj}$だから
\begin{equation*}
	\Tr{(AB)}= \sum_{j= 1}^n(BA)_{jj}= \Tr{(BA)}
\end{equation*}

	\item 
\begin{itemize}
	\item $\Tr{(ABC)}= \Tr{(CAB)}$\\
	$\Tr{(A(BC))}= \Tr{((BC)A)}$より$\Tr{(ABC)}= \Tr{(CAB)}$
	\item $\Tr{(BCA)}= \Tr{(CAB)}$\\
	$\Tr{(B(CA))}= \Tr{((CA)B)}$より$\Tr{(BCA)}= \Tr{(CAB)}$
\end{itemize}

	\item $\Tr{(P^{-1}(AP))}= \Tr{((AP)P^{-1})}$
	行列の結合法則より$(AP)P^{-1}= A(P^{-1}P)= AE= A$。したがって$\Tr{P^{-1}AP}= \Tr{A}$

	\item $A= \begin{pmatrix}1& 0\\ 0& 0\end{pmatrix},\ B= \begin{pmatrix}0& 1\\ 0& 0\end{pmatrix},\ C= \begin{pmatrix}0& 0\\ 1& 0\end{pmatrix}$とする。
\begin{itemize}
	\item $ABC$
\begin{equation*}
	ABC= \begin{pmatrix}1& 0\\ 0& 0\end{pmatrix}\begin{pmatrix}0& 1\\ 0& 0\end{pmatrix}\begin{pmatrix}0& 0\\ 1& 0\end{pmatrix}= \begin{pmatrix}1& 0\\ 0& 0\end{pmatrix}
\end{equation*}
	$\Tr{(ABC)}= 1$

	\item $ACB$
\begin{equation*}
	ACB= \begin{pmatrix}1& 0\\ 0& 0\end{pmatrix}\begin{pmatrix}0& 0\\ 1& 0\end{pmatrix}\begin{pmatrix}0& 1\\ 0& 0\end{pmatrix}= \begin{pmatrix}0& 0\\ 0& 0\end{pmatrix}
\end{equation*}
	$\Tr{(ACB)}= 0$
\end{itemize}
	よって$\Tr{(ABC)}\ne \Tr{(ACB)}$
\end{enumerate}
\end{souanswer}

\begin{tcolorbox} %%R6_3rd_2
	$\fbox{2}$ $A= \begin{pmatrix}1& i\\ {-i}& 1\end{pmatrix}$とする。ただし$i$は虚数単位とする。
\begin{enumerate}
	\item $A$の固有値を複素数の範囲ですべて求めよ。
	\item (1)で求めた$A$の各固有値の固有ベクトルを求めよ。
	\item $U^{-1}AU$が対角行列となるようなユニタリ行列$U$を求めよ。
	\item $n$を自然数とするとき、$A^n$を求めよ。
\end{enumerate}
\end{tcolorbox}

\begin{souanswer} %%R6_3rd_2_ans
	$\fbox{2}$ 
\begin{enumerate}
	\item 固有値を$\lambda$とする。固有多項式は
\begin{align*}
	\det{(A- \lambda E)}
&= \begin{vmatrix}
	1- \lambda& i\\
	-i& 1- \lambda
\end{vmatrix}\\
&= (1- \lambda)^2- i\cdot (-i)\\
&= \lambda(\lambda- 2)
\end{align*}
	よって求める固有値は$\lambda= 0, 2$
	\item 
\begin{enumerate}
	\item $\lambda= 2$のときの固有ベクトルを$\bm{v}_1=\begin{pmatrix}x_1\\ y_1\end{pmatrix}$とする。
\begin{equation*}
	\begin{pmatrix}-1& i\\ -i& -1\end{pmatrix}\begin{pmatrix}x_1\\ y_1\end{pmatrix}= \begin{pmatrix}0\\ 0\end{pmatrix}
\end{equation*}
	$y_1= c_1\ (c_1\colon 任意定数)$とすると、$x_1= ic_1$.
	よって$\lambda= 0$のときの固有ベクトルは$\bm{v}_1= c_1\begin{pmatrix}i\\ 1\end{pmatrix}$

	\item $\lambda= 0$のときの固有ベクトルを$\bm{v}_2= \begin{pmatrix}x_2\\ y_2\end{pmatrix}$とする。
\begin{equation*}
	\begin{pmatrix}1& i\\ -i& 1\end{pmatrix}\begin{pmatrix}x_2\\ y_2\end{pmatrix}= \begin{pmatrix}0\\ 0\end{pmatrix}
\end{equation*}
	$y_2= c_2 (c_2\colon 任意定数)$とすると、$x_2= -ic_2$.
	よって$\lambda= 2$のときの固有ベクトルは$\bm{v}_2= c_2\begin{pmatrix}-i\\ 1\end{pmatrix}$
\end{enumerate}
	\item 固有ベクトル$\bm{v}_1, \bm{v}_2$を正規化する。また正規化したベクトルをそれぞれ$\bm{u}_1, \bm{u}_2$とかく。
\begin{align*}
	\bm{u}_1&= \frac{1}{\sqrt{i^2+ 1^2}}\begin{pmatrix}i\\ 1\end{pmatrix}= \frac{1}{\sqrt{2}}\begin{pmatrix}i\\ 1\end{pmatrix}\\
	\bm{u}_2&= \frac{1}{\sqrt{(-i)^2+ 1^2}}\begin{pmatrix}-i\\ 1\end{pmatrix}= \frac{1}{\sqrt{2}}\begin{pmatrix}-i\\ 1\end{pmatrix}
\end{align*}
	対角行列は$\begin{pmatrix}2& 0\\ 0& 0\end{pmatrix}$であるから
\begin{equation*}
	U^{-1}AU= U^*AU= \begin{pmatrix}2& 0\\0& 0\end{pmatrix}
\end{equation*}
	これをみたすような$U$は$\frac{1}{\sqrt{2}}\begin{pmatrix}i& -i\\1& 1\end{pmatrix}$
	\item $A^n= UD^nU^{-1}$が成り立つ。ここで$D$は対角行列
\begin{align*}
	A^n
	&= \frac{1}{2}\begin{pmatrix}i& -i\\ 1& 1\end{pmatrix}\begin{pmatrix}2^n& 0\\ 0& 0\end{pmatrix}\begin{pmatrix}-i& 1\\ i& 1\end{pmatrix}\\
	&= \frac{1}{2}\begin{pmatrix}2^n& i\cdot 2^n\\ -i\cdot 2^n& 2^n\end{pmatrix}\\
	&= 2^{n- 1}A
\end{align*}
\end{enumerate}
\end{souanswer}

\begin{tcolorbox} %%R6_3rd_3
	$\fbox{3}$ つぎの問いに答えよ。ただし$e$を自然対数の底とする。
\begin{enumerate}
	\item 実数上で定義された実数値関数$f(x)$が$x= a$において微分可能であることの定義を$\epsilon- \delta$論法を用いて書け。
	\item 実数上の関数$f(x)= (\sin{x}+ 2)^x$を微分せよ。
	\item 実数$x$、正の実数$y, z$に対して定義される3変数関数$g(x, y, z)= ex(z^y- y^z)- e^x$を考える。
\begin{enumerate_roman}
	\item $g(x, y, z)= 0$で定義される曲面の$(x, y, z)= (1, 1, 2)$における接平面の方程式を求めよ。
	\item $g(x, y, z)= 0$で定義される陰関数$z= z(x, y)$の1階偏導関数$z_x, z_y$を求めよ。また2階偏導関数$z_{yy}$の$(x, y)= (1, 1)$での値$z_{yy}(1, 1)$を求めよ。
\end{enumerate_roman}
\end{enumerate}
\end{tcolorbox}

\begin{souanswer} %%R6_3rd_3_ans
	$\fbox{3}$
\begin{enumerate}
	\item 任意の$\epsilon> 0$に対して、ある$\delta> 0$が存在し、$	0< |x- a|< \delta$ならば、$0< |f(x)- f(a)|< \epsilon$
\begin{tcolorbox}[title= 論理記号ver.]
	$\forall \epsilon>0, \exists \delta> 0, 0< |x- a|< \delta \Longrightarrow 0< |f(x)- f(a)|< \epsilon$
\end{tcolorbox}

	\item 両辺の対数をとる。ここで$\sin{x}+ 2> 0$であるから真数条件をみたす。
\begin{align*}
	\log{f(x)}&= x\log{(\sin{x}+ 2)}\\
	\frac{f'(x)}{f(x)}&= \log{(\sin{x}+ 2)}+ x\cdot \frac{\cos{x}}{\sin{x}+ 2}\\
	f'(x)&= \left(\log{(\sin{x}+ 2)}+ \frac{x\cos{x}}{\sin{x}+ 2}\right)\cdot(\sin{x}+ 2)^x
\end{align*}

\begin{enumerate_roman}
	\item 求める接平面の方程式は
\begin{equation} \tag{P}\label{eq:求める接平面の方程式}
	g_x(1, 1, 2)(x- 1)+ g_y(1, 1, 2)(y- 1)+ g_z(1, 1, 2)(z- 2)= 0 
\end{equation}
	それぞれの偏微分の値を計算する。
\begin{align*}
	g_x(1, 1, 2)
	&= e(z^y- y^z)- e^x\\
	&= e(2^1- 1^2)- e^1\\
	&= 0\\
	g_y(1, 1, 2)
	&= ex(z^y\log{z}- zy^{z- 1})\\
	&= e\cdot s(2^1\log{2}- 2\cdot 1^{2- 1})\\
	&= 2e(\log{2}- 1)\\
	g_z(1, 1, 2)
	&= ex(yz^{y- 1}- y^z\log{y})\\
	&= e\cdot 1(2\cdot 1^{2- 1}- 1^{2- 1}\log{1})\\
	&= e
\end{align*}
	\eqref{eq:求める接平面の方程式}に代入して
\begin{equation*} %なんか変になってる気がする
	2e(\log{2}- 1)(y- 1)+ e(z- 2)= 0
\end{equation*}
	$e$で割って
\begin{equation*}
	(\log{2}- 1)y+ z- 2\log{2}= 0
\end{equation*}
	\item 陰関数の微分法より$z_x= -\frac{g_x}{g_z}, z_y= -\frac{g_y}{g_z}$となる。
\begin{align*}
	z_x&= \frac{e^x- e(z^y- y^z)}{ex(yz^{y- 1}- y^z\log{y})}\\
	z_y&= -\frac{ex(z^y\log{z}- zy^{z- 1})}{ex(yz^{y- 1}- y^z\log{y})}\\
	&= -\frac{z^y\log{z}- zy^{z- 1}}{yz^{y- 1}- y^z\log{y}}
\end{align*}
	次に$z_{yy}(1, 1)$を求める。条件より$(x, y)= (1, 1)$のとき$z= 2$であり、$z_y(1, 1)$は先ほどの偏微分係数を用いて求められる。
\begin{equation*}
	z_y(1, 1)= -\frac{g_y(1, 1, 2)}{g_z(1, 1, 2)}= -\frac{2e(\log{2}- 1)}{e}= -2(\log{2}- 1)= 2(1- \log{2})
\end{equation*}
	商の微分と連鎖律を用いて$z_{yy}= \frac{\partial}{\partial y}\left(-\frac{g_y}{g_z}\right)$を計算する。
\begin{equation} \label{eq:zyy}
	z_{yy}= -\frac{(g_{yy}+ g_{yz}z_y)g_z- g_y(g_{zy}+ g_{zz}z_y)}{{g_z}^2}
\end{equation}
	各2階偏微分導関数を計算し、$(1, 1, 2)$の値を定める。
\begin{itemize}
	\item $g_{yy}= ex(z^y(z^y(\log{z})^2)- z(z- 1)y^{z- 2})$
	\item[$\Rightarrow$] $g_{yy}(1, 1, 2)= 2e((\log{2})^2- 1)$ 
	\item $g_{yz}= g_{zy}= ex(yz^{y- 1}\log{z}+ z^{y- 1}- y^{z- 1}- zy^{y- 1}\log{y})$
	\item[$\Rightarrow$] $g_{yz}(1, 1, 2)= e\log{2}$
	\item $g_{zz}= ex(y(y- 1)z^{y- 2}- y^z(\log{y})^2)$
	\item[$\Rightarrow$] $g_{zz}(1, 1, 2)= 0$
\end{itemize}
	これらを\eqref{eq:zyy}に代入して
\begin{align*}
	z_{yy}
	&= -\frac{2e^2(\log{2}- 1)- 2e^2\log{2}(\log{2}- 1)}{e^2}\\
	&= -2(\log{2}- 1)+ 2\log{2}(\log{2}- 1)\\
	&= -2(\log{2}- 1)(1- \log{2})\\
	z_{yy}(1, 1)
	&= 2(\log{2}- 1)^2
\end{align*}



\end{enumerate_roman}
\end{enumerate}
\end{souanswer}

\begin{tcolorbox} %%R6_3rd_4
	$\fbox{4}$ 次の問いに答えよ。
\begin{enumerate}
	\item 極座標表示$x= r\cos{\theta}, y= r\sin{\theta}$に対して、Jacobian$\frac{\partial(x, y)}{\partial(r, \theta)}$を求めよ。
	\item $\int x^2\cdot 2^{x^3} dx$を求めよ。
	\item $\iint_{D} 2^{(x^2+ y^2)^{3/2}} dxdy$を求めよ。ただし
\begin{equation*}
	D= \{(r\cos{\theta}, r\sin{\theta});\ 0\le \theta\le 1,\ 2\theta\le r\le 2\}
\end{equation*}
\end{enumerate}
\end{tcolorbox}

\begin{souanswer} %%R6_3rd_4_ans
	$\fbox{4}$
\begin{enumerate}
	\item 
\begin{equation*}
	\frac{\partial (x, y)}{\partial (r, \theta)}= \begin{vmatrix}{\partial x}/{\partial r}& {\partial y}/{\partial r}\\ {\partial x}/{\partial \theta}& {\partial y}/{\partial \theta}\end{vmatrix}
\end{equation*}
	それぞれの偏微分の値を求める。
	${\partial x}/{\partial r}= \cos{\theta}$, 
	${\partial y}/{\partial r}= \sin{\theta}$, 
	${\partial x}/{\partial \theta}= -r\sin{\theta}$, 
	${\partial y}/{\partial \theta}= r\cos{\theta}$.\\
	求めるJacobianは$\begin{vmatrix}\cos{\theta}& \sin{\theta}\\ -r\sin{\theta}& r\cos{\theta}\end{vmatrix}= r$

	\item $t= x^3$とおく。$dt= 3x^2 dx$
\begin{align*}
	I
	&= \int 2^t \left(\frac{1}{3}\right) dt\\
	&= \frac{1}{3\log{2}} 2^{x^3}+ C\ \ (C:積分定数)
\end{align*}

	\item 座標変換して$x^2+ y^2= r^2$、(1)よりJacobianは$r$であるから$dxdy= r drd\theta$
\begin{equation*}
	I= \int_{\theta= 0}^1 \int_{r= 2\theta}^2 2^{r^3}\cdot rdrd\theta
\end{equation*}

\begin{tcolorbox}
\begin{itemize}
	\item $r$の範囲: $r$の最小値は$2\theta= 0(\leadsto \theta= 0)$、最大値は2。範囲は$0\le \theta\le 2$
	\item $\theta$の範囲: $2\theta\le r$より$\theta\le \frac{r}{2}$、$0\le \theta$。範囲は$0\le \theta\le \frac{r}{2}$
\end{itemize}
\end{tcolorbox}
\begin{align*}
	I
	&= \int_{r= 0}^2\int_{\theta= 0}^{r/2} r\cdot 2^{r^3}drd\theta\\
	&= \int_{r= 0}^2\left[\theta\cdot r\cdot 2^{r^3}\right]_{\theta= 0}^{r/2}dr\\
	&= \frac{1}{2}\int_{r= 0}^2 r^2\cdot 2^{r^3}dr\\
	&= \frac{1}{2}\left[\frac{1}{3\log{2}}\cdot 2^{r^3}\right]_{r= 0}^2\\
	&= \frac{1}{6\log{2}}\cdot (2^8- 2^0)\\
	&= \frac{85}{2\log{2}}
\end{align*}
\end{enumerate}
\end{souanswer}



\chapter{令和7年度}
\section{一般第1期}
\begin{tcolorbox} %%R7_1st_1
	$\fbox{1}$ $A= \begin{pmatrix}1& 1& 1\\ 1& 0& 1\\ {-1}& {-1}& {-1}\end{pmatrix}$とし、$X$を$n$次複素正方行列で、ある非負整数$k$に対して$X^k$が零行列となるとする。
\begin{enumerate}
	\item $A$の固有多項式・固有値を求めよ。
	\item $X$の固有値をすべて求めよ。
	\item $X^m$が零行列となる最小の非負整数を$m$とすると、$m\le n$であることを示せ。
	\item $X$が対角化可能ならば、$X$は零行列であることを示せ。
\end{enumerate}
\end{tcolorbox}

\begin{souanswer} %%R7_1st_1_ans
	$\fbox{1}$
\begin{enumerate}
	\item 固有値を$\lambda$とする。
\begin{align*}
	\begin{vmatrix}A- \lambda E\end{vmatrix}
&= 
\begin{vmatrix}
	1- \lambda& 1& 1\\
	1& -\lambda& 1\\
	-1& -1& -1- \lambda
\end{vmatrix}\\
&= \lambda(1- \lambda)(1+ \lambda)- \lambda\\
&= -\lambda^3
\end{align*}
\begin{tabular}{lll}
	固有方程式& $\colon$& $\begin{vmatrix}A- \lambda E\end{vmatrix}= -\lambda^3$\\
	固有値& $\colon$& 0
\end{tabular}

	\item 固有値を$\mu$、対応する固有ベクトルを$\bm{v}\ (\bm{v}\ne \bm{0})$とする。固有方程式より
\begin{equation*}
	X\bm{v} = \mu \bm{v}
\end{equation*}
	この式の両辺に左から繰り返し$X$をかけ続ける。ここで任意の非負整数$k$に対して$X^k= O$であったから
\begin{equation*}
	X^k\bm{v}= \mu^k \bm{v}\iff Ov= \mu^k \bm{v}\iff \bm{0}= \mu^k \bm{v}
\end{equation*}
	ここで$\bm{v}= \bm{0}$であるから$\mu^k= 0$となるしかない。$\mu$を複素数の範囲で解くと$\mu= 0$
\end{enumerate}
\end{souanswer}


\begin{tcolorbox} %%R7_1st_2
	$\fbox{2}$ $n$次正方行列$H$が次の条件を満たすとする。
\begin{enumerate_alpha}
	\item 各要素が1または-1
	\item 各行が直交する。すなわち$h_i$を$H$の第$i$行の行ベクトルとすると、$j\ne k$ならば$h_j{}^t h_k= 0$ただし、行列$M$に対しては${}^t M$はその転置行列とする。
\end{enumerate_alpha}
	次の問に答えよ。
\begin{enumerate}
	\item $H^t H= nI_n$を示せ。ただし$I_n$は$n\times n$の単位行列とする。
	\item  $H$の逆行列を求めよ。
	\item $H$の各列も直交することを示せ。
	\item $n$を3以上の奇数とすると、条件(a), (b)をともに満たす$n\times n$行列は存在しないことを示せ。
\end{enumerate}
\end{tcolorbox}

\begin{souanswer} %%R7_1st_2_ans
	$\fbox{2}$
\begin{enumerate}
	\item 
	\item 
	\item 
	\item 
\end{enumerate}
\end{souanswer}

\begin{tcolorbox} %%R7_1st_3
	$\fbox{3}$ 以下の問いに答えよ．
	$g_1(x, y)= (e^y+ e^{-y})\cos{x}$, $g_2(x, y)= -(e^y- e^{-y})\sin{x}$、および$f(x, y)= \exp{(g_1(x, y))}\cos{(g_2(x, y))}$と定義する。
\begin{enumerate}
	\item $\frac{\partial f}{\partial x}$および$\frac{\partial f}{\partial y}$を求めよ。
	\item $\frac{\partial^2 f}{\partial x^2}+ \frac{\partial^2 f}{\partial y^2}$を求めよ。
	\item 次を満たす$x, y$の実数値関数$h(x, y)$が存在する場合にはその関数を求めよ。また存在しない場合にはその証明を書け。
\begin{equation*}
	\frac{\partial f}{\partial x}(x, y)= \frac{\partial h}{\partial y}(x, y) および\frac{\partial f}{\partial y}(x, y)= -\frac{\partial h}{\partial x}(x, y)
\end{equation*}
\end{enumerate}
\end{tcolorbox}

\begin{souanswer} %%R7_1st_3_ans
	$\fbox{3}$
\begin{enumerate}
	\item 
	\item 
	\item 
\end{enumerate}
\end{souanswer}


\begin{tcolorbox} %%R7_1st_4
	$\fbox{4}$ $a, b$を正の定数とする。3次元空間$\R^3$において、$z= x^2+ 2y^2$を$S$として$(x_0, y_0, z_0)$を$S$上の点とする。
\begin{enumerate}
	\item\label{vol:楕円柱} 曲面$S$、平面$z= 0$、楕円柱$\frac{x^2}{a^2}+ \frac{y^2}{b^2}= 1$で囲まれた立体の体積を求めよ。
	\item $a^2 +2b^2= 1$において\ref{vol:楕円柱}の立体の体積の最大値を求めよ。
	\item 曲面$S$の円柱$x^2+ 4y^2= 1$で囲まれた部分の曲面積を求めよ。
	\item 曲面$S$の$(x_0, y_0, z_0)$における接平面の方程式を求めよ。
\end{enumerate}
\end{tcolorbox}

\begin{souanswer} %%R7_1st_4_ans
	$\fbox{4}$
\begin{enumerate}
	\item 
	\item 
	\item 
	\item 
\end{enumerate}
\end{souanswer}



\section{一般第2期}
\begin{tcolorbox} %%R_7_2nd_1
	$\fbox{1}$
\begin{enumerate}
	\item $V, W$を体$K$上のベクトル空間とする。写像$f\colon V\to W$が$K$上の線形写像であることの定義を述べよ。
	\item $\C[x]$を複素数体$\C$に係数をもち、$x$を変数とする一変数多項式の集合とする。$\C[x]$は通常の多項式の和と定数倍によりベクトル空間をなす。$n$次多項式$\bm{v}= a_n x^n+ a_{n- 1}x^{n- 1}+ \cdots+ a_1 x+ a_0\ (n> 0ならばa_n\ne 0)$に対し、$f(\bm{v})$を次のように定めた写像$f\colon \C[x]\to \C[x]$が線形写像であるかを理由とともに答えよ。
\begin{enumerate_alpha}
	\item $f(\bm{v})= a_n+ a_{n- 1}+ \cdots+ a_0$
	\item $f(\bm{v})= a_n\times a_{n- 1}\times \cdots a_0$
	\item $f(\bm{v})= a_{n- 1} x^n+ a_{n- 2} x^{n- 1}+ \cdots+ a_0 x+ a_n$
\end{enumerate_alpha}
	\item $n$を正整数とする。$M_n(\C)$を複素体$\C$の要素を成分に持つ$n$次正方行列の集合とする。行列の和と定数倍によりベクトル空間をなす。つぎの写像が線形写像であるかを判定せよ。
\begin{enumerate_alpha}
	\item トレース$\Tr\colon M_n(\C)\to \C$
	\item 行列式$\det\colon M_n(\C)\to \C$
\end{enumerate_alpha}
\end{enumerate}
\end{tcolorbox}

\begin{souanswer} %%R7_2nd_1_ans
	$\fbox{1}$
\begin{enumerate}
	\item ベクトル空間$V, W$との間の写像$f\colon V\to W$が線型写像であるとは、任意のベクトル$\bm{u}, \bm{v}\in V$とスカラー
$k\in \R$に対して
\begin{enumerate_roman}
	\item\label{law:加法性} 加法性：$f(\bm{u}+ \bm{v})= f(\bm{u})+ f(\bm{v})$
	\item\label{law:斉次性} 斉次性：$f(k\bm{v})= kf(\bm{v})$
\end{enumerate_roman}
	\item 多項式$v= \sum_{i= 0}^n a_ix^i$に対して判定する。
\begin{enumerate_alpha}
	\item 線型写像。多項式に$x= 1$を代入した$f(v)= v(1)$を返す写像。任意の多項式$u, v$とスカラー$k\in \R$に対して
\begin{align*}
	(u+ v)(1)&= u(1)+ v(1)\\
	(kv)(1)&= kv(1)
\end{align*}
	が成り立つため線型写像。

	\item 線型写像ではない。斉次性を満たさない。$\bm{v}= x+ 1(a_1= 1, a_0= 1)$とし、スカラー$k= 2$とすると、
\begin{itemize}
	\item $f(2\bm{v})= f(2x+ 2)= 2\times 2= 4$
	\item $2f(\bm{v})= 2(1\times 1)= 2\times 1= 2$
\end{itemize}
	となり、線型写像ではない。

	\item 線型写像ではない。加法性を満たさない。$\bm{u}= x(n= 1, a_1= 1, a_0= 1)\Rightarrow f(\bm{v})= 0\cdot x^1+ 1= 1$と$\bm{v}= 1(n= 0, a_1= 1\Rightarrow f(\bm{v})= 1)$を考える。
\begin{itemize}
	\item $f(\bm{u})+ f(\bm{v})= 1+ 1= 2$
	\item $\bm{u}+ \bm{v}= 1\cdot x^1+ 1= x+ 1$
\end{itemize}
	したがって$f(\bm{u}+ \bm{v})\ne f(\bm{u})+ f(\bm{v})$
\end{enumerate_alpha}

	\item 
\begin{enumerate_alpha}
	\item トレース$\Tr{}\colon M_n(\C)\to \C$は線型写像。
	定義より$\Tr{A}= \sum_{i= 1}n a_{ii}$
\begin{itemize}
	\item $\Tr{(A+ B)}= \sum_{i= 1}^n(a_{ii}+ b_{ii})= \sum_{i= 1}^n a_{ii}+ \sum_{i= 1}^n b_{ii}= \Tr{A}+ \Tr{B}$
	\item $\Tr{(kA)}= \sum_{i= 1}^n (ka_{ii})=k\sum_{i= 1}^n a_{ii}= k\Tr{A}$
\end{itemize}


	\item 行列式$\det{}\colon M_n(\C)\to \C$は線型写像ではない。一般に$|A+ B|\ne |A||B|$。また斉次性についても$|kA|= k^n|A|$となり、$n> 1$では一致しない。
\end{enumerate_alpha}
\end{enumerate}
\end{souanswer}

\begin{tcolorbox} %R7_2nd_2
	$\fbox{2}$ $n$次実正方行列$A$が次の条件をみたすとき、$A$を正定値行列という。ただし$\tr{\bm{v}}$は$\bm{v}$の転置とする。
\begin{tcolorbox}
	零ベクトルでない任意の$n$次元実数ベクトル$\bm{v}$に対し、$\tr{\bm{v}} A\bm{v}> 0$
\end{tcolorbox}
	$H_n$を$i$行$j$列の要素が$\frac{1}{i+ j- 1}$である$n$次実正方行列とする。
\begin{enumerate}
	\item $H_3$の階数($\rank{}$)を求めよ。
	\item $H_2$は定数値行列であることを示せ。
	\item $n$次元ベクトル$\bm{v}= \tr{(v_1\ v_2\cdots\ v_n)}$に対し、以下を示せ。
\begin{equation*}
	\tr{\bm{v}} H_n\bm{v}= \sum_{i= 1}^n\sum_{j= 1}^n \frac{v_i v_j}{i+ j-1}
\end{equation*}
	\item $n$次元ベクトル$\bm{v}= \tr{(v_1\ v_2\cdots\ v_n)}$に対し、多項式$f_v(x)$を次のように定める
\begin{equation*}
	f_v(x)= v_n x^{n-1}+ v_{n- 1} x^{n- 2}+ \cdots+ v_2 x+ v_1
\end{equation*}
	以下を示せ。
\begin{equation*}
	\tr{\bm{v}} H_n\bm{v}= \int_0^1 {f_v(x)}^2 dx
\end{equation*}
	\item $H_n$が正定値行列であることを示せ。
\end{enumerate}
\end{tcolorbox}

\begin{souanswer} %%R7_2nd_2_ans
	$\fbox{2}$
\begin{enumerate}
	\item $H_3$は以下のように書き下せる。
\begin{equation*}
	H_3= 
\begin{pmatrix}
	1& 1/2& 1/3\\
	1/2& 1/3& 1/4\\
	1/3& 1/4& 1/5
\end{pmatrix}
\end{equation*}
	行列式を調べると、
\begin{equation*}
	\begin{vmatrix}H_3\end{vmatrix}= 1\cdot \left(\frac{1}{15}- \frac{1}{16}\right)- \frac{1}{2}\left(\frac{1}{10}- \frac{1}{12}\right)+ \frac{1}{3}\left(\frac{1}{8}- \frac{1}{9}\right)= \frac{1}{240}- \frac{1}{120}+ \frac{1}{216}= \frac{1}{21600}\ne 0
\end{equation*}
	$\begin{vmatrix}H_3\end{vmatrix}\ne 0$なのでフルランク。$\rank{H_3}= 3$

	\item $H_2= \begin{pmatrix}1& 1/2\\ 1/2& 1/3\end{pmatrix}$で任意の非零ベクトル$\bm{v}= \tr{\begin{pmatrix}v_1& v_2\end{pmatrix}}$に対して
\begin{align*}
	\tr{\bm{v}}H_2\bm{v}
	&= \begin{pmatrix}v_1\\ v_2\end{pmatrix}\begin{pmatrix}1& 1/2\\ 1/2& 1/3\end{pmatrix}\begin{pmatrix}v_1& v_2\end{pmatrix}\\
	&= v_1^2+ v_1v_2+ \frac{1}{3}v_2^2\\
	&= \left(v_1+ \frac{1}{2}v_2\right)^2+ \frac{1}{12}v_2^2
\end{align*}
	となる。$v_1, v_2$のうち少なくとも一方が0でなければ、この値は常に正となる。すなわち$H_2$は正定値行列。

	\item $\bm{v}= (v_1, \ldots, v_n)^\mathrm{T}$、および$H_n$の要素$(H_n)_{ij}= \frac{1}{i+ j- 1}$より
\begin{align*}
	\bm{v}^\mathrm{T}H_n\bm{v}
	&= \sum_{i= 1}^n v_i (H_n)_i\\
	&= \sum_{i= 1}^n v_i\left(\sum_{j= 1}^n \frac{1}{i+ j- 1}v_j\right)\\
	&= \sum_{i= 1}^n \sum_{j= 1}^n\frac{v_i v_j}{i+ j- 1}
\end{align*}

	\item $f_v(x)= \sum_{k= 1}^n v_kx^{k- 1}$とおくと
\begin{align*}
	\{f_v(x)\}^2
	&= \left(\sum_{i= 1}^n v_ix^{i- 1}\right)\left(\sum_{j= 1}^n v_jx^{j- 1}\right)\\
	&= \sum_{i= 1}^n\sum_{j= 1}^n v_iv_j x^{i+ j- 2}
\end{align*}
	これを0から1まで積分する。
\begin{equation*}
	\int_0^1\{f_v(x)\}^2 dx= \sum_{i= 1}^n \sum_{j= 1}^n v_iv_j \int_0^1 x^{i+ j-2} dx
\end{equation*}
	ここで$\int_0^1 x^{i+ j- 2} dx= \frac{1}{i+ j- 1}$なので、
\begin{equation*}
	\int_0^1\{f_v(x)\}^2 dx= \sum_{i= 1}^n \sum_{j= 1}^n\frac{v_iv_j}{i+ j- 1}
\end{equation*}

	\item 任意の$\bm{v}\ne 0$に対して、$\bm{v}^\mathrm{T}H\bm{v}= \int_0^1 \{f_v(x)\}^2 dx$
\begin{itemize}
	\item $\{f_v(x)\}^2\ge 0$であるからその積分値は0以上
	\item  $\int_0^1 \{f_v(x)\}^2 dx= 0$となるのは連続関数$f_v(x)$が区間$[0, 1]$で恒等的に0である時に限られる。
	\item $f_v(x)$は$n- 1$次多項式であり、$\bm{v}\ne 0$であれば、係数の少なくとも1つが0ではないため、$f_v(x)$は恒等的に0にはならない。
\end{itemize}


\end{enumerate}
\end{souanswer}

\begin{tcolorbox} %%R7_2nd_3
	$\fbox{3}$
\begin{enumerate}
	\item $n$次元実数空間$\R^n$における有界閉集合$D$上で定義された連続関数は最大値および最小値をとることを示せ。なお有界閉集合$D$上の任意の点列は$D$に収束する部分列をもつことを用いてよい。
	\item $s+ t= 1,\ 0< s, t< 1$および$d> 1$とする。
\begin{equation*}
	s\log{(sd)}+ t\log{(td)}- s\log{(d- 1)}\ge 0
\end{equation*}
	であることを示せ。

	\item $a> b> c> 0,\ d> 3$とする。また3次元実数空間$\R^3$上の関数$f(x, y, z),\ \phi(x, y, z)$を
\begin{align*}
	f(x, y, z)= x^2+ y^2+ z^2,
	\phi(x, y, z)= e^{ax^2}+ e^{by^2}+ e^{cz^2}- d
\end{align*}
	とする。$\phi(x, y, z)= 0$のもとで$f(x, y, z)$の最大値および最小値を求めよ。
\end{enumerate}
\end{tcolorbox}

\begin{souanswer} %%R4_3nd_3_ans
	$\fbox{3}$
\begin{enumerate}
	\item 
\begin{enumerate_alpha}
	\item \textbf{有界性} 
	$f(D)$が上に有界であることを示す。有界でないならば任意の$k\in \N$に対し、$f(x_k)> k$となる$x_k\in D$が存在する。
	$D$は有界閉集合なので点列$\{x_k\}$は$D$内の点$x_0$に収束する部分列$\{x_{k_j}\}_{j= 1}^\infty$が存在する。$f$は連続であるから$\lim_{j\to \infty} x_{k_j}= f(x_0)$となるが、これは$f(x_{k_j})> k_j\to \infty$と矛盾。したがって$f(D)$は有界。

	\item \textbf{最大値の存在} 
	$M= \sup_{x\in D}f(x)$とおく。上限の定義より$f(x_k)\to M$となり、点列$\{x_{k_j}\}\subset D$が選べる。同様に$x_0\in D$に収束する部分列$\{x_{k_j}\}$をとると連続性より$f(x_0)= \lim_{j\to \infty} f(x_{k_j})= M$となり、$f$は点$x_0$で最大値を取る。最小値についても同様に示せる。
% もう少し丁寧に
\end{enumerate_alpha}

	\item $g(s):= s\log{(sd)}+ t\log{(td)}- s\log{(d- 1)}$とおく。$t= 1- s$なので
\begin{equation*}
	g(s)= s\log{(sd)}+ (1- s)\log{((1- s)d)}- s\log{(d- 1)}
\end{equation*}
	$s$で微分すると
\begin{align*}
	g'(s)
	&= \log{(sd)}+ 1- \log{((1- s)d)}- 1- \log{(d- 1)}\\
	&= \log{\left(\frac{s}{(1- s)(1- d)}\right)}
\end{align*}
\begin{equation*}
	g'(s)= 0\iff \frac{s}{(1- s)(1- d)}= 1\iff s= \frac{d- 1}{d}
\end{equation*}
	この前後で$g'(s)$の符号が変化するため$s= \frac{d- 1}{d}$で最小値をとる。
\begin{equation*}
	g\left(\frac{d- 1}{d}\right)= \frac{d- 1}{d}\log{(d- 1)}+ \frac{1}{d}\log{1}- \frac{d- 1}{d}\log{(d- 1)}= 0
\end{equation*}
	よって$g(s)\ge 0$が示された。

	\item Lagrangeの未定乗数法を用いる。
\begin{equation*}
	L(x, y, z; \lambda):= x^2+ y^2+ z^2- \lambda(e^{ax^2}+ e^{by^2}+ e^{cz^2}- d)
\end{equation*}
	とおき、偏微分が0となる点を求める。
\begin{itemize}
	\item ${\partial L}/{\partial x}= 2x- \lambda(2axe^{ax^2})= 2x(1- \lambda ae^{ax^2})= 0$
	\item ${\partial L}/{\partial y}= 2y- \lambda(2bye^{by^2})= 2y(1- \lambda be^{by^2})= 0$
	\item ${\partial L}/{\partial z}= 2z- \lambda(2cze^{cz^2})= 2z(1- \lambda ce^{cz^2})= 0$
\end{itemize}
	各変数について座標が0または$e^{k\xi^2}= \frac{1}{\lambda k}$となる必要がある。% どゆこと？
\begin{itemize}
	\item \textbf{最大値} $a> b> c> d$なのでもっとも広がりやすいのは係数が小さい$z$方向。もし$x= y= 0$となる$L(0, 0, 2)= 1+ 1+ e^{cz^2}-d= 0$より$e^{cz^2}= d- 2$
	このとき$f= z^2= \frac{1}{c}\log{(d- 2)}$。$d> 3$より$\log{(d- 2)}> 0$で成立。
	同様の候補は$\frac{1}{a}\log{(d- 2)}, \frac{1}{b}\log{(d- 2)}$であるが、係数が最小の$c$のとき$f$は最大になる。

	\item \textbf{最小値} すべての変数が0でない場合を考える。$1= \lambda ae^{ax^2}= \lambda be^{by^2}= \lambda ce^{cz^2}$とすると、$e^{ax^2}= 1/\lambda a, e^{by^2}= 1/\lambda b, e^{cz^2}= 1/\lambda c$を制約式に代入して$\frac{1}{\lambda a}+ \frac{1}{\lambda b}+ \frac{1}{\lambda} c= d\Rightarrow \frac{1}{\lambda}= \frac{d}{\frac{1}{a}+ \frac{1}{b}+ \frac{1}{c}}$
	これを$x^2= \frac{1}{a}\log{\left(\frac{1}{\lambda a}\right)}$等に代入して足し合わせる。境界や対称性を考慮すると、係数が最大の軸上に点が乗るときが最も原点に近くなる。
\end{itemize}
\end{enumerate}
\end{souanswer}

\begin{tcolorbox} %%R7_2nd_4
	$\fbox{4}$
\begin{enumerate}
	\item $f(x)= \cos{x}$の$n$階導関数を求め、$x= 0$でのTaylor展開$\sum_{n= 0}^\infty \frac{f^{(n)}(0)}{n!}x^n$を求めよ。
	\item\label{eq:g(x)の定義} $
	g(x)= 
\begin{cases}
\begin{matrix}
	e^{-\frac{1}{x^2}}& x\ne 0\\
	0& x= 0
\end{matrix}
\end{cases}
$
	とする。$g'(x)$を求めよ。
	\item $n$を自然数としたとき\eqref{eq:g(x)の定義}で定義された$g(x)$の$n$階導関数が$2n- 2$次以下の多項式$p_n(x)$を用いて、
$
	g^{(n)}(x)=
\begin{cases}
\begin{matrix}
	\frac{p_n(x)}{x^{3n}}e^{-\frac{1}{x^2}}& x\ne 0\\
	0& x= 0
\end{matrix}
\end{cases}
$
	となることを示せ。
	\item \eqref{eq:g(x)の定義}で定義された$g(x)$は$x= 0$でTaylor展開できないことを証明せよ。
\end{enumerate}
\end{tcolorbox}

\begin{souanswer} %%R7_2nd_4_ans
	$\fbox{4}$
\begin{enumerate}
	\item $\cos{x}$の$n$次導関数は以下のように書ける。
\begin{equation*}
	f^{(n)}(x)= \cos{\left(x+ \frac{n\pi}{2}\right)}
\end{equation*}
	次に$x= 0$の値を求める。
\begin{itemize}
	\item $n= 2k(偶数)$のとき：
\begin{equation*}
	f^{(2k)}(0)= \cos{k\pi}= (-1)^k
\end{equation*}
 
	\item $n= 2k+ 1(奇数)$のとき：
\begin{equation*}
	f^{(2k+ 1)}(0)= \cos{\left(k\pi+ \frac{\pi}{2}\right)}= 0
\end{equation*}
\end{itemize}
	したがってTalyor展開は以下の通り
\begin{equation*}
	\sum_{n= 0}^\infty \frac{f^{(n)}}{n!}x^n= \sum_{n= 0}^\infty \frac{(-1)^k}{(2k)!}x^{2k}
\end{equation*}

	\item $x\ne 0$のときは合成関数の微分を用いる。
\begin{align*}
	g'(x)
	&= e^{-\frac{1}{x^2}}\cdot \frac{d}{dx}(-x^{-2})\\
	&= e^{-\frac{1}{x^2}}\cdot (2x^{-3})\\
	&= \frac{2}{x^3}e^{-\frac{1}{x^2}}
\end{align*}
	$x= 0$における微分係数は微分の定義より求める。
\begin{equation*}
	g'(0)= \lim_{h\to 0}\frac{g(h)- g(0)}{h}= \lim_{h\to 0} \frac{e^{-\frac{1}{h^2}}}{h}
\end{equation*}
	ここで$t= \frac{1}{h}$とおくと$h\to 0$のとき$|t|\to \infty$となり、
\begin{equation*}
	g'(0)= \lim_{|t|\to \infty} \frac{t}{e^{t^2}}= 0
\end{equation*}
	よってすべての$x$で
\begin{equation*}
	g'(x)= 
\begin{cases}
	\frac{2}{x^3}e^{-\frac{1}{x^2}}\ \ (x\ne 0)\\
0\ \ (x= 0)
\end{cases}
\end{equation*}
	\item\label{step:R7-2nd数学的帰納法} 数学的帰納法で示す。
\begin{enumerate_roman}
	\item \textbf{$n= 1$のとき} 
	$g^{(1)}= \frac{2}{x^3}e^{-\frac{1}{x^2}}$。$p_1(x)= 2$とすれば$2(1)- 2= 0$次の多項式であり、成立。

	\item \textbf{$n= k$の成立を仮定し、$x\ne 0$で計算} 
\begin{align*}
	g^{(k+ 1)}(x)
	&= \frac{d}{dx}\left[p_k(x)x^{-3k}e^{-x^{-2}}\right]\\
	&= \left[p'_k(x)x^{-3k}- 3kp_k(x)x^{-3k- 1}+ p_k(x)x^{-3k}(2x^{-3})\right]e^{-x^{-2}}\\
\end{align*}
	　分子を$p_{k+ 1}(x)$とおくと、$p_k(x)$が$2k- 2$次なので$p_{k+ 1}(x)$の最高時は$2k- 2+ 2= 2(k+ 1)- 2$次となり、次数条件をみたす。

	\item \textbf{$x= 0$で証明}
	$g^{(k+ 1)}(0)= 0$を仮定すると
\begin{align*}
	g^{(k+ 1)(0)}
	&= \lim_{x\to 0} \frac{g^{(k)}(x)- 0}{x}\\
	&= \lim_{x\to 0} \frac{p_k(x)}{x^{3k+ 1}}e^{-\frac{1}{x^2}}
\end{align*}
	ここで$t= 1/x$とすると$x\to 0$のとき$|t|\to \infty$
\begin{equation*}
	\lim_{|t|\to 0} \frac{t^{3k+ 1}p_k(1/t)}{e^{t^2}}= 0
\end{equation*}
	となり、成立する。
\end{enumerate_roman}
	\item 関数$f(x)$が$x= 0$でTaylor展開できるとは$x= 0$の近傍で
\begin{equation*}
	g(x)= \sum_{n= 0}^\infty \frac{g^{(n)}(0)}{n!}x^n
\end{equation*}
	が成立することである。しかし\eqref{step:R7-2nd数学的帰納法}より$g^{(n)}(0)= 0$であることが示されている。もしTaylor展開可能なら
\begin{equation*}
	g(x)= \sum_{n= 0}^\infty \frac{0}{n!}x^n= 0
\end{equation*}
	となり、常に右辺が収束する。一方定義より$x\ne 0$において$g(x)= e^{-1/x^2}> 0$。したがって$x= 0$のどのような近傍でも
\begin{equation*}
	g(x)\ne \sum_{n= 0}^\infty \frac{g^{(n)}(0)}{n!}x^n
\end{equation*}
	となり、$g(x)$は$x= 0$でTaylor展開できない。
\end{enumerate}
\end{souanswer}

\section{一般第3期}
	志願者不在のため実施無し


\chapter{令和8年度}
\section{一般第1期}

\section{一般第2期}

\section{一般第3期}
\begin{tcolorbox} %%R8_3rd_1
	$\fbox{1}$
\begin{enumerate}
	\item $A_3= \begin{pmatrix}1& 2& 0\\ 0& 1& 2\\ 2& 0& 1\end{pmatrix}$とする。
\begin{enumerate_alpha}
	\item $A_3$の固有値を複素数の範囲ですべて求めよ。
	\item $A_3$の固有値のうち、絶対値が最大のものの固有ベクトルを複素数の範囲ですべて求めなさい。
\end{enumerate_alpha}
	\item $n$を2以上の整数とする。$n$次正方行列$B_n$は、$i$行$j$成分$B_n(i, j)$が以下で与えられているとする。
\begin{equation*}
	B_n({i, j})= 
\begin{cases}
	i+ j- 1\ \ (i+ j\le n)\\
	i+ j- n- 1\ \ (i+ j\ge n+ 1) 
\end{cases}
\end{equation*}
	とする。
\begin{enumerate_alpha}
	\item $B_4, B_5$を行列で表しなさい。
	\item $B_n$が固有値に$\frac{n(n- 1)}{2}$をもつことを示せ。
	\item $B_n$の固有値の中で絶対値が最大なものは$\frac{n(n- 1)}{2}$であることを示せ。
\end{enumerate_alpha}
\end{enumerate}
\end{tcolorbox}

\begin{souanswer} %R8_3rd_1_ans

\end{souanswer}

\begin{tcolorbox} %%R8_3rd_2
	$\fbox{2}$ $A$を複素正方行列とし、$A^*$は$A$の随伴行列$\tr{\bar{A}}$とする。$A= A^*$が成り立つとき$A$をエルミート行列、$A= -A^*$を満たすとき、$A$を歪エルミート行列という。
\begin{enumerate}
	\item 以下の行列に対してエルミート行列，歪エルミート行列、以外を理由とともに判定をせよ
\begin{enumerate_alpha}
	\item $\begin{pmatrix}1& i& -i&\\ -i& 1& 0\\ i& 0& 1\end{pmatrix}$
	\item $\begin{pmatrix}i& -1& -1\\ 1& i& 0\\ 1& 0& i\end{pmatrix}$
	\item $\begin{pmatrix}1+ i& 1& 1\\ 1& i& 0\\ 1& 0& i\end{pmatrix}$
\end{enumerate_alpha}
	\item 歪エルミート行列の対角成分は、0または純虚数であることを示せ。
	\item エルミート行列かつ歪エルミート行列は存在するか。
	\item $M$を複素正方行列とする。$M- M^*$がエルミート行列であることを示せ。
	\item $n$を正整数とする。$n$次正方行列$M, N$がエルミート行列であるとき、$i(MN- NM)$がエルミート行列であることを示せ。
	\item 歪エルミート行列の固有値は0または純虚数であることを示せ。
\end{enumerate}
\end{tcolorbox}

\begin{souanswer} %R8_3rd_2_ans
\begin{enumerate}
	\item 
\begin{enumerate_alpha}
	\item 
\begin{equation*}
	A^*= \tr{\begin{pmatrix}1& -i& i\\ i& 1& 0\\ -i& 0& 1\end{pmatrix}}= \begin{pmatrix}1& i& -i\\ -i& 1& 0\\ i& 0& 1\end{pmatrix}= A
\end{equation*}	
	$A^*= A$となるためエルミート行列。
	\item 
\begin{equation*}
	A^*= \tr{\begin{pmatrix}-i& -1& -1\\ 1& -1& 0\\1& 0& -i\end{pmatrix}}= \begin{pmatrix}-i& 1& 1\\ -1& -1& 0\\ -1& 0& -i\end{pmatrix}= -A
\end{equation*}
	$A^*= -A$となるため歪エルミート行列。
	\item 
\begin{equation*}
	\tr{\begin{pmatrix}1+ i& 1& 1\\ 1& i& 0\\1& 0& i\end{pmatrix}}= \begin{pmatrix}1+ i& 1& 1\\ 1& i& 0\\ 1& 0& i\end{pmatrix}
\end{equation*}
	$A^*\ne A$かつ$A^*\ne -A$であるためエルミート行列でも歪エルミート行列でもない。
\end{enumerate_alpha}
	\item 行列$A$を歪エルミート行列とする。すなわち$A^*= -A$。$A$の$(k, k)$対角成分を$a_{kk}$とかく。対角成分に注目すると$\bar{\tr{A}}$の対角成分は変化しないため
\begin{equation} \label{R8_3rd_2(2)_diagonal_reration}
	\bar{a_{kk}}= -a_{kk}
\end{equation}
	が成り立つ。ここで$a_{kk}= x+ yi\ (x, y\in \R)$とすると
\begin{equation*}
	\bar{a_{kk}}= x- yi
\end{equation*}
	\zcref{R8_3rd_2(2)_digonal_reration}に代入すると
\begin{equation*}
	x- yi= -(x- yi)\Leftrightarrow 2x= 0\Leftrightarrow x= 0
\end{equation*}
	よって$a_{kk}= yi$となる。$y\ne 0$のときは純虚数となるため、対角成分は0または純虚数となる。
	\item ある行列$A$がエルミート行列かつ歪エルミート行列であると仮定する。すなわち$A^*= A$かつ$A^*= -A$。よって$A= -A$。これをみたすのは$A= O(零行列)$のみ。
	\item 行列$X:= M- M^*$とし、$X^*$を計算する。
\begin{equation*}
	X*= (M- M^*)^*= M^*- M^{**}= M^* -M= -(M- M^*)= -X\ (\because 随伴行列の線形性)
\end{equation*}
	$X^*= -X$となったから$M- M^*$は歪エルミート行列
	\item 仮定より$M= M^*$かつ$N= N^*$。行列$X:= i(MN- NM)$に対して$X^*= X$であることを示す。随伴行列の性質より
\begin{equation*}
	X^*= (i(MN- NM))^*= -i(MN- NM)^*= -i(N^*M^*- M^*N^*)= -i(NM- MN)= X
\end{equation*}
	$X^*= X$となるため$i(MN- NM)$はエルミート行列。
	\item 行列$A$を歪エルミート行列とし、$\lambda$をその固有値、$\bm{x}$を対応する固有ベクトルとする。すなわち$A\bm{x}= \lambda\bm{x}$。
	この式の両辺の左から$\bm{x}^*$をかけて
\begin{equation} \label{R8_3rd_2(5)_reration}
	\bm{x}^*A\bm{x}= \bm{x}^*\lambda\bm{x}
\end{equation}
	次に$\bm{x}^*A\bm{x}$の全体の随伴をとると
\begin{equation*}
	(\bm{x}^*A\bm{x})^*= \bm{x}^*A^*(\bm{x}^*)^*= \bm{x}^*A^*\bm{x}
\end{equation*}
	$A$は歪エルミートであるから$A^*= -A$を代入して
\begin{equation*}
	\bm{x}^*A^*\bm{x}= -\bm{x}^*A\bm{x}
\end{equation*}
	ここで\zcref{R8_3rd_2(5)_reration}を代入すると
\begin{equation*}
	(\lambda\bm{x}^*\bm{x})^*= -(\lambda\bm{x}^*\bm{x})
\end{equation*}
	スカラーの共役の性質と$\bm{x}^*\bm{x}= \norm{x}^2$であることを用いると
\begin{equation*}
	\bar{\lambda}\norm{\bm{x}}^2= -\lambda\norm{x}^2
\end{equation*}
	$\bm{x}$は固有ベクトルであるから$\norm{\bm{x}}\ne 0$。両辺を$\norm{\bm{x}}$で割ると
\begin{equation*}
	\bar{\lambda}= -\lambda
\end{equation*}
	$\lambda= a+ bi\ (a, b\in \R)$とすると
\begin{equation*}
	a- bi= -(a- bi)\Leftrightarrow 2a= 0
\end{equation*}
	よって$\lambda$は0または純虚数であることが示された。
\end{enumerate}
\end{souanswer}

\begin{tcolorbox} %%R8_3rd_3
	$\fbox{3}$
\begin{enumerate}
	\item $x\ne 0$のとき$\frac{d}{dx}\left(\frac{\cos{x}}{x}\right)$を求めよ。
	\item $\int_1^\infty \frac{1}{x^2}dx$を求めよ。
	\item $\left|\int_1^\infty \frac{\cos{x}}{x^2}dx\right|\le 1$を求めよ。
	\item $\left|\int_1^\infty \frac{\sin{x}}{x} dx\right|< \infty$を示せ。
	\item $\int_1^\infty \left|\frac{\sin{x}}{x}\right|dx= \infty$を示せ。	
\end{enumerate}
\end{tcolorbox}

\begin{souanswer} %R8_3rd_3_ans
\begin{enumerate}
    \item\label{R8_3rd_3(1)} 商の微分公式より
\begin{equation*}
    \frac{d}{dx}\left(\frac{\cos{x}}{x}\right)= \frac{-\sin{x}\cdot x- \cos{x}\cdot 1}{x^2}= \frac{-x\sin{x}- \cos{x}}{x^2}
\end{equation*}
    \item\label{R8_3rd_3(2)} 広義積分の定義より
\begin{equation*}
    \int_1^\infty \frac{1}{x^2}dz
	= \lim_{R\to \infty}\int_z^R \frac{1}{x^2}dx
    = \lim_{R\to \infty} \left[-\frac{1}{x}\right]_1^R
    = \lim_{R\to \infty}\left(-\frac{1}{R}- (-1)\right)= 0+ 1= 1
\end{equation*}
    \item\label{R8_3rd_3(3)} 絶対値の積分に関する三角不等式$\left|\int_a^b f(x)dx\right|\le \int_a^b|f(x)|dx$を用いる。
\begin{equation*}
	\left|\int_1^\infty \frac{\cos{x}}{x^2}dx \right|
	\le \int_1^\infty \left|\frac{\cos{x}}{x^2}\right|dx
	= \int_1^\infty \frac{|\cos{x}|}{x^2}dx
\end{equation*}
	すべての実数$x$に対して$|\cos{x}|\le 1$であることと積分範囲$x\ge 1$において$x^2> 0$であるから、
\begin{equation*}
	\int_1^\infty \frac{|\cos{x}|}{x^2}dx\le \int_1^\infty \frac{1}{x^2}dx
\end{equation*}
	\zcref{R8_3rd_3(2)}の結果より$\int_1^\infty \frac{1}{x^2}dx= 1$であるから、
\begin{equation*}
	\left|\int_1^\infty \frac{\cos{x}}{x^2}dx\right|\le 1
\end{equation*}
	\item \zcref{R8_3rd_3(1)}の両辺を1から$\infty$まで積分すると、
\begin{equation*}
	\int_a^\infty \frac{\sin{x}}{x}dx= -\int_1^\infty \frac{\cos{x}}{x^2}dx- \left[\frac{\cos{x}}{x}\right]_1^\infty
\end{equation*}
	ここで第2項の極限は、
\begin{equation*}
	\left[\frac{\cos{x}}{x}\right]_1^\infty= \lim_{R\to \infty}\frac{\cos{R}}{R}= \frac{\cos{1}}{1}= -\cos{1}
\end{equation*}
	したがって
\begin{equation*}
	\int_a^\infty \frac{\sin{x}}{x}dx= -\int_1^\infty \frac{\cos{x}}{x^2}dx+ \cos{1}
\end{equation*}
	両辺の絶対値をとり、三角不等式を適用すると、
\begin{equation*}
	\left|\int_a^\infty \frac{\sin{x}}{x}dx\right|\le \left|-\int_1^\infty \frac{\cos{x}}{x^2}dx\right|+ |\cos{1}|
\end{equation*}
	\zcref{R8_3rd_3(3)}の結果より$\left|\int_1^\infty \frac{\cos{x}}{x}dx\right|\le 1$であり、$\cos{1}$は有限の値であるから
\begin{equation*}
	\left|\int_1^\infty \frac{\sin{x}}{x}dx\right|\le 1+ |\cos{x}|< \infty
\end{equation*}
	\item 積分区間を$\pi$ごとに分割して、下から評価する。任意の自然数$N$に対して、積分範囲を拡大しながら考える。
\begin{equation*}
	\int_1^{N\pi} \left|\frac{\sin{x}}{x}\right|dx\ge \int_\pi^{N\pi}|\left|\frac{\sin{x}}{x}\right|dx= \sum_{k= 1}^{N- 1}\int_{k\pi}^{(k+ 1)\pi} \frac{|\sin{x}|}{x}dx
\end{equation*}
	各区間$x\in [k\pi, (k+ 1)\pi]$において$x\le (k+ 1)\pi$であるため$\frac{1}{x}\ge \frac{1}{(k+ 1)\pi}$が成り立つ。これを適用すると、
\begin{equation*}
	\sum_{k= 1}^{N- 1}\int_{k\pi}^{(k+ 1)\pi}\frac{|\sin{x}|}{x}dx\ge \sum_{k= 1}^{N- 1} \int_{k\pi}^{(k+ 1)\pi}|\sin{x}|dx
\end{equation*}
	ここで半周期の正弦波の面積は常に2であるから
\begin{equation*}
	\sum_{k= 1}^{N- 1}\frac{2}{(k+ 1)\pi}= \frac{2}{\pi}\sum_{k= 1}^{N- 1} \frac{1}{k+ 1}= \frac{2}{\pi}\left(\frac{1}{2}+ \frac{1}{3}+ \cdots+ \frac{1}{N}\right)
\end{equation*}
	$N\to \infty$とすると調和級数$\sum \frac{1}{k}$は$\infty$に発散するから、下から評価した級数が$\infty$に発散するならばもとの積分も発散する。
\begin{equation*}
	\int_1^\infty \left|\frac{\sin{x}}{x}\right|dx= \infty
\end{equation*}
\end{enumerate}
\end{souanswer}

\begin{tcolorbox} %%R8_3rd_4
	$\fbox{4}$
\begin{enumerate}
	\item $\frac{d}{dr}\exp{(-r^2)}$を求めなさい。
	\item $x= r\cos{\theta}, y= r\sin{\theta}$とする。ヤコビアン${\partial(x, y)}/{\partial(r, \theta)}$を求めよ。
	\item $\int_0^\infty \int_0^\infty \exp{(-x^2- y^2)}dxdy$を求めよ。
	\item $\int_0^\infty \int_0^\infty\int_0^\infty \exp{(-x^2- y^2- z^2)}dxdydz$を求めよ。
\end{enumerate}
\end{tcolorbox}

\begin{souanswer} %R8_3rd_4_ans
\begin{enumerate}
	\item 合成関数の微分法を用いて
\begin{equation*}
	\frac{d}{dx}\exp(-r^2)= \exp(-r^2)\cdot \frac{d}{dr}(-r^2)= -2r\exp{(-r^2)}
\end{equation*}
	\item\label{R8_3rd_4(2)} $x= r\cos{\theta}, y= r\sin{\theta}$と変数変換すると、
\begin{equation*}
	\det{\begin{pmatrix}dx/dr& dx/d\theta\\ dy/dr& dyd/\theta\end{pmatrix}}= r
\end{equation*}	
	\item $x= r\cos{\theta}, y= r\sin{\theta}$と変数変換すると、
\begin{itemize}
	\item $\exp{(-x^2- y^2)}= \exp{(-r^2(\cos^2\theta+ \sin^2\theta))}= \exp{-r^2}$
	\item 積分領域は$x\ge 0, y\ge 0$だから$0\le r\le \infty, 0\le \theta\le 2\pi$
	\item ヤコビアンは\zcref{R8_3rd_4(2)}より$r$だから$dxdy= rdrd\theta$
\end{itemize}
	\item \label{R8_3rd_4(4)}
\begin{align*}
	\int_0^\infty \int_0^\infty \exp{(-x^2- y^2)}dxdy
	&= \int_{\theta= 0}^{2\pi}\int_{r= 0}^\infty \exp{(-r^2)}rdrd\theta\\
	&= \int_{r= 0}^\infty \exp{(-r^2)}rdr \int_{\theta= 0}^{2\pi}d\theta\\
	&= \left[-\frac{1}{2}\exp{(-r^2)}\right]_0^\infty \left[\theta\right]_{\theta= 0}^{2\pi}\\
	&= \frac{1}{2}\cdot \frac{\pi}{2}= \frac{\pi}{4}
\end{align*}
	\item 積分は以下のように書き換えることができる。
\begin{align*}
	\int_0^\infty \int_0^\infty\int_0^\infty \exp{(-x^2- y^2- z^2)}dxdydz
	&= \left(\int_0^\infty \exp(-x^2)dx\right)	\left(\int_0^\infty \exp(-y^2)dy\right)\left(\int_0^\infty \exp(-z^2)dz\right)\\
	&= I^3
\end{align*}
	\zcref{R8_3rd_4(4)}より$I> 0$であるから$I= \frac{\sqrt{\pi}}{2}$。よって求める積分の値は$I^3= \left(\frac{\sqrt{\pi}}{2}\right)^3= \frac{\pi\sqrt{\pi}}{8}$
\end{enumerate}
\end{souanswer}













\end{document}